% Auth: Nicklas Vraa
% Docs: https://github.com/NicklasVraa/LiX
% Everything you need to know about this template is found in on the github repository above. Stars are very appreciated.

\documentclass{novel}

% ===== PACKAGES POUR LE CONTENU SCIENTIFIQUE =====
% Mathématiques avancées
\usepackage{amsmath, amssymb, amsthm}
\usepackage{mathtools}
\usepackage{bm} % pour les vecteurs en gras
\usepackage{mdframed}

% Figures et graphiques
\usepackage{graphicx}
\usepackage{subcaption}
\usepackage{wrapfig}
\usepackage{float}

% Tableaux améliorés
\usepackage{array}
\usepackage{booktabs}
\usepackage{multirow}
\usepackage{longtable}
\usepackage{pdfpages}

% ===== PACKAGES POUR LA MODÉLISATION =====
% Algorithmes et pseudo-code
\usepackage{algorithm}
\usepackage{algorithmic}

% Diagrammes et schémas
\usepackage{tikz}
\usetikzlibrary{arrows, positioning, shapes}

% Listings de code (si nécessaire)
\usepackage{listings}
\usepackage{xcolor}

% Numérotation des équations par chapitre
\numberwithin{equation}{chapter}

% ===== COMMANDES PERSONNALISÉES =====

% Environnements théoriques
\theoremstyle{definition}
\newtheorem{definition}{Définition}[chapter]
\newtheorem{exemple}{Exemple}[chapter]
\newtheorem{notation}{Notation}[chapter]

\theoremstyle{plain}
\newtheorem{theoreme}{Théorème}[chapter]
\newtheorem{proposition}{Proposition}[chapter]
\newtheorem{lemme}{Lemme}[chapter]
\newtheorem{proprietes}{Propriété}[chapter]
\newtheorem{corollaire}{Corollaire}[chapter]
\newtheorem{exercice}{Exercice}[chapter]

\theoremstyle{remark}
\newtheorem{remarque}{Remarque}[chapter]
\newtheorem{objectif}{Objectif}[chapter]
%\newtheorem{note}{Note}[chapter]

\theoremstyle{proof}
\newtheorem{solution}{Solution}[chapter]


\lang      {french}
\title     {Estimation}
\subtitle  {Fondements et méthodes d'estimation en statistique inférentielle}
\authors   {Lorenzo Segoni}
\cover     {resources/novel_front.pdf}{resources/novel_back.pdf}
%\license   {CC}{by-nc-sa}{3.0}
%\isbn      {978-0201529838}
%\publisher {NV Publishing Company}
%\edition   {1}{2023}
%\dedicate  {Myself}{Because I'm cool}
%\thank     {Thank you to me for being the best}
%\keywords  {fiction, template, packages}

%\note{Lorem ipsum dolor sit amet, consectetur adipiscing elit. Praesent porttitor est arcu, sed euismod metus imperdiet iaculis. Quisque vestibulum molestie nulla, non consectetur tellus mollis a. Nunc commodo magna a elit commodo dignissim.}

\blurb{Ce cours présente les principes fondamentaux de la modélisation statistique et de l'estimation de paramètres inconnus à partir de données échantillonnées. Il introduit rigoureusement les propriétés essentielles que l'on attend d'un estimateur, telles que l'absence de biais, la consistance asymptotique et la minimisation du risque quadratique. Le document détaille ensuite les approches majeures de construction d'estimateurs, notamment la méthode des moments et le maximum de vraisemblance, tout en explorant les limites théoriques de leur précision via l'information de Fisher et la borne de Cramér-Rao. Enfin, le cours expose la démarche de quantification de l'incertitude à travers l'élaboration d'intervalles de confiance, qu'ils soient exacts grâce à la méthode des quantités pivotales dans les modèles gaussiens, ou asymptotiques. L'ensemble fournit ainsi un cadre mathématique solide pour interpréter des observations empiriques et maîtriser l'aléa statistique.}

\begin{document}

\toc


\chapter{Méthodes d'estimation}

\section{Exemple introductif : Exemple du sondage}

Un projet de loi est soumis à referendum. La proportion $\theta$ d'individu qui vont voter "oui" est un réel inconnu de $[0,1]$. Pour savoir si le oui va l'emporter, on réalise un sondage et on interroge $n$ individus dans la population française. On tire ces individus au hasard et avec remplacement.

On note $X_i = 1$ si la réponse du $i$-éme individu indépendantes et identiquement distribuées, de loi de Bernoulli $\mathcal{B}(\theta)$. Au vu des réalisations des variables aléatoires $X_i$, on cherche à déterminer la valeur de $\theta$ et savoir si le projet de loi va être adopté : on souhaite \textit{estimer} $\theta$ et \textit{tester} si $\theta$ est supérieur ou non à $\frac{1}{2}$. Il est naturel de considérer comme estimateur de $\theta$ la moyenne empirique:

$$
\hat{\theta} = \frac{1}{n} \sum_{i=1}^n X_i
$$  

L'inégalité de Tchebychev nous dit que, pour tout $\varepsilon >0$,

$$
\mathbb{P}_\theta \left( |\hat{\theta} - \theta | \geq \varepsilon \right) \leq \frac{\mathbb{V}[X_i]}{n \varepsilon ^2} \leq \frac{1}{4 n \varepsilon^2}
$$

Noter que la loi de $\theta$ dépend du $\theta$ inconnu, d'où l'indice dans la probabilité.

Puisque la variance d'une loi de Bernoulli est $\theta (1-\theta)$ et est majorée uniformément par $\frac{1}{4}$ le choix de $\varepsilon = \frac{1}{2 \sqrt{n \alpha}}$ conduit à :

$$
\mathbb{P}_\theta \left( |\hat{\theta} - \theta | \leq \frac{1}{2 \sqrt{n \alpha}} \right) \geq 1 - \alpha
$$

On dit que $[\hat{\theta} - \frac{1}{2 \sqrt{n \alpha}} ; \hat{\theta}  + \frac{1}{2 \sqrt{n \alpha}}]$ est un intervelle de condiance pour $\theta$ de niveeau de confiance $1 - \alpha$.

Pour $\alpha = 5 \%$ et $n=1000$, la demi-longuer de l'intervalle est environ égale à $0.0707$.

Ce résultat montre un biais fréquent dans l'interprétation des sondages par les médias. Par exemple, si un journal annonce une variation de $2 \%$ des intentions de vote au cours de $6$ mois sur la base d'un échantillon de $1000$ personnes, ce changement n'a aucune valeur. Comme on a montré, ce mouvement apparent n'est que du bruit statistique.

\newpage

\section{Modèles statistiques}

\begin{definition}
On appelle \textbf{observation} la variable aléatoire $X = ( X_1 , \cdots , X_n)$

On dit que $( X_1 , \cdots , X_n)$  est un $n$-\textbf{échantillon} de loi $P$ si les variables aléatoires $X_1 , \cdots , X_n$ sont indépendantes identiquement distrubuées (i.i.d.) de loi $P$.
\end{definition}

\smallskip

\begin{mdframed}
\begin{definition}[Modèle statistique]
Un modèle statistique est un triplet :
\[
(E, \mathcal{E}, (P_\theta)_{\theta \in \Theta})
\]
où :
\begin{itemize}
  \item $(E, \mathcal{E})$ est un espace mesurable probabilisable (dans ce cours, on retiendra surtout que $X \in E$) ;
  \item $(P_\theta)_{\theta \in \Theta}$ est une famille de lois de probabilité sur $E$ ;
  \item $\Theta$ est l'ensemble des paramètres.
\end{itemize}

Si $\Theta \subset \mathbb{R}^d$ avec $d \geq 1$, le modèle est dit \textbf{paramétrique}.  
Sinon, il est dit \textbf{non paramétrique}.
\end{definition}
\end{mdframed}

\smallskip

\begin{remarque}
Un modèle est toujours faux : la réalité est trop complexe pour être décrite parfaitement. L'enjeu est que l'approximation apportée par le modèle soit suffisamment fidèle pour répondre à la question posée.  

En pratique, ce sont souvent les spécialistes du domaine (médecine, biologie, économie, etc.) qui proposent ou discutent le modèle.
\end{remarque}

\begin{exemple}
\begin{enumerate}
  \item \textbf{Jeu de pile ou face répété $n$ fois :}
  \[
  E = \{0,1\}^n, \quad 
  \Theta = \,]0,1[\,, \quad
  P_\theta = \big(\mathcal{B}(\theta)\big)^{\otimes n}.
  \]
  Une observation est $X = (X_1,\ldots,X_n) \in \{0,1\}^n$, avec $X_i$ des variables aléatoires indépendantes de loi $\mathcal{B}(\theta)$.

  \item \textbf{Modélisation de la taille d'individus par une loi normale.} On observe sur un échantillon de taille $n$ :
  \[
  E = \mathbb{R}^n, \quad X = (X_1,\ldots,X_n),
  \]
  où les $X_i$ sont des variables aléatoires indépendantes de loi $\mathcal{N}(m,\sigma^2)$, avec
  \[
  \Theta = \{ (m,\sigma^2) : m \in \mathbb{R}, \ \sigma^2 > 0 \} \subset \mathbb{R}^2.
  \]

  \item \textbf{Même population, mais modélisation par une loi à densité sur $[0.5, 2.5]$ :}
  \[
  \Theta = \{ f : [0.5,2.5] \to \mathbb{R}_+ \ \mid f \text{ est une densité} \}.
  \]
  Il s'agit alors d'un \textbf{modèle non paramétrique}.
\end{enumerate}
\end{exemple}

\paragraph{Valeur vraie du paramètre.}
Le point clé de la modélisation est qu'il existe une valeur \emph{vraie} $\theta^* \in \Theta$ telle que :
\[
P_X = P_{\theta^*}.
\]
Mais $\theta^*$ est inconnue.

Dans toute la suite, on notera simplement $\theta$ pour $\theta^*$.

On utilisera les notations :
\[
\mathbb{P}_\theta \quad \text{et} \quad \mathbb{E}_\theta
\]
pour bien marquer la dépendance en $\theta$. On a ainsi :
\[
P(X \in A) = P_X(A) = P_{\theta^*}(A), \qquad P_\theta(X \in A) = P_\theta(A).
\]

\begin{remarque}
On considérera souvent une observation $X$ comme un $n$-échantillon de variables aléatoires i.i.d.\ de loi $Q_\theta$ :
\[
P_\theta = (Q_\theta)^{\otimes n}.
\]
On peut alors définir le modèle statistique comme la famille $\{Q_\theta, \theta \in \Theta\}$.
\end{remarque}

\begin{definition}[Identifiabilité]
Un modèle $\{P_\theta, \theta \in \Theta\}$ est dit \textbf{identifiable} si :
\[
P_{\theta_1} = P_{\theta_2} \quad \Rightarrow \quad \theta_1 = \theta_2.
\]
\end{definition}

\begin{remarque}
L'identifiabilité est une condition \emph{nécessaire} pour pouvoir estimer $\theta$ de manière cohérente : si deux paramètres distincts produisent la même loi, aucune observation ne permet de les distinguer. Dans toute la suite, on supposera le modèle identifiable.

\medskip

\textbf{Exemple.} Considérons $X \sim \mathcal{N}(\mu, \sigma^2)$ avec $\theta = (\mu, \sigma^2)$. La loi de $X$ détermine de façon unique $(\mu, \sigma^2)$ : ce modèle est identifiable. En revanche, si l'on paramétrisait par $(\mu, \sigma)$ et $(\mu, -\sigma)$, les deux valeurs du paramètre donneraient la même loi, et le modèle ne serait plus identifiable.
\end{remarque}

\subsection{Estimateur statistique}

On suppose que le modèle $\{P_\theta, \theta \in \Theta\}$ est identifiable.  
On observe $X$ et on s'intéresse à une quantité d'intérêt $g(\theta)$, où :
\[
g : \Theta \to \mathbb{R}^p, \quad p \geq 1.
\]

\begin{mdframed}
\begin{definition}[Statistique et estimateur]
\begin{itemize}
  \item Une \textbf{statistique} $T(X)$ est une fonction mesurable de $X$ (éventuellement de paramètres connus), qui ne dépend pas de $\theta$.
  \item Un \textbf{estimateur} de $g(\theta)$ est une statistique $\hat{g} = h(X)$ à valeurs dans $\mathbb{R}^p$ (avec $h$ mesurable).
\end{itemize}
\end{definition}
\end{mdframed}

\begin{exemple}
Pour un $n$-échantillon $(X_1,\ldots,X_n)$, les objets suivants sont des statistiques :
\[
\overline{X}_n, \quad X_1, \quad X_2^2 \cdot X_8, \quad 0.
\]
En revanche, $\overline{X}_n + \theta$ \emph{n'est pas} une statistique car elle dépend du paramètre inconnu $\theta$.
\end{exemple}

\begin{remarque}
Si l'on dispose d'une réalisation $x$ de $X$, alors $\hat{g}(x)$ est appelé un \textbf{estimé} de $g(\theta)$.  
Par abus, on note souvent $\hat{g}$ à la fois l'estimateur (variable aléatoire) et son estimé (valeur numérique).
\end{remarque}

\begin{notation}[Estimateur plug-in]
Si $\hat{\theta}$ est un estimateur de $\theta$, un estimateur naturel de $g(\theta)$ est :
\[
g(\hat{\theta}).
\]
On appelle cela un \textbf{estimateur plug-in}.  
Sa qualité dépend de la qualité de $\hat{\theta}$ et de la régularité de $g$.
\end{notation}

\begin{definition}[Risque quadratique]
Le risque quadratique de $\hat{g}$ est défini, pour tout $\theta \in \Theta$, par :
\[
R(\hat{g}, g(\theta)) = \mathbb{E}_\theta \big[ (\hat{g} - g(\theta))^2 \big] \in [0, +\infty].
\]
\end{definition}

\begin{remarque}
\textbf{Lien avec l'inégalité de Markov :}
Par l'inégalité de Markov,
\[
\mathbb{P}_\theta \big( |\hat{g} - g(\theta)| > \varepsilon \big) \leq \frac{R(\hat{g}, g(\theta))}{\varepsilon^2}.
\]
Ainsi, plus le risque quadratique est petit, plus l'estimateur est concentré autour de la valeur vraie.
\end{remarque}

\paragraph{Existe-t-il un estimateur optimal ?}

On pourrait souhaiter qu'il existe un estimateur $\hat{g}$ meilleur que tous les autres pour tout $\theta$, c'est-à-dire :
\[
\exists\, \hat{g} \text{ tel que } \forall\, \theta \in \Theta,\ \forall\, \hat{g}' \text{ estimateur},\quad R(\hat{g}, g(\theta)) \leq R(\hat{g}', g(\theta)).
\]
La réponse est en général \textbf{non}.

\begin{exemple}[Comparaison de deux estimateurs]
Considérons le modèle :
\[
\{ (\mathcal{N}(\theta,1))^{\otimes n}, \ \theta \in \mathbb{R} \}.
\]

Deux estimateurs de $\theta$ :
\[
\hat{\theta}_1 = \overline{X}_n, \qquad \hat{\theta}_2 = 0.
\]

On a :
\[
R(\hat{\theta}_1, \theta) = \frac{1}{n}, \qquad R(\hat{\theta}_2, \theta) = \theta^2.
\]

De plus :
\[
\mathbb{E}_\theta[\hat{\theta}_1] = \theta, \qquad \mathrm{Var}_\theta(\hat{\theta}_1) = \frac{1}{n}.
\]

Conclusion :
\begin{itemize}
  \item $\hat{\theta}_1$ est sans biais avec un risque constant $1/n$.
  \item $\hat{\theta}_2$ est biaisé, mais pour $|\theta| < \frac{1}{\sqrt{n}}$, son risque est inférieur à celui de $\hat{\theta}_1$.
\end{itemize}

Ainsi, même un estimateur « stupide » comme $\hat{\theta}_2$ peut être meilleur que l'estimateur usuel $\overline{X}_n$ pour certaines valeurs de $\theta$. Bien sûr, plus $n$ est grand, plus l'intervalle de $\theta$ où $\hat{\theta}_2$ est meilleur devient petit.
\end{exemple}

\begin{mdframed}
\begin{proposition}[Décomposition biais-variance]
Pour tout $\theta \in \Theta$, on a :
\[
R(\hat{g}, g(\theta)) 
= \big( \mathbb{E}_\theta[\hat{g}] - g(\theta) \big)^2 
+ \mathbb{E}_\theta\!\left[ \big(\hat{g} - \mathbb{E}_\theta[\hat{g}]\big)^2 \right].
\]
Autrement dit :
\[
R(\hat{g}, g(\theta)) = B(\hat{g}, g(\theta))^2 + \mathrm{Var}_\theta(\hat{g}),
\]
où $B(\hat{g}, g(\theta))$ désigne le biais de $\hat{g}$.
\end{proposition}
\end{mdframed}

\begin{definition}[Biais]
Le biais d'un estimateur $\hat{g}$ est défini par :
\[
\theta \mapsto B(\hat{g}, g(\theta)) = \mathbb{E}_\theta[\hat{g}] - g(\theta).
\]
On dit que $\hat{g}$ est \textbf{sans biais} si, pour tout $\theta \in \Theta$ :
\[
\mathbb{E}_\theta[\hat{g}] = g(\theta).
\]
\end{definition}

\begin{proof}
On part de l'identité :
\[
\hat{g} - g(\theta)
= \big( \hat{g} - \mathbb{E}_\theta[\hat{g}] \big)
+ \big( \mathbb{E}_\theta[\hat{g}] - g(\theta) \big).
\]
En élevant au carré :
\[
(\hat{g} - g(\theta))^2
= \big(\hat{g} - \mathbb{E}_\theta[\hat{g}]\big)^2
+ 2\big(\hat{g} - \mathbb{E}_\theta[\hat{g}]\big)\big(\mathbb{E}_\theta[\hat{g}] - g(\theta)\big)
+ \big(\mathbb{E}_\theta[\hat{g}] - g(\theta)\big)^2.
\]
En prenant l'espérance sous $\mathbb{P}_\theta$, le terme croisé disparaît car
\[
\mathbb{E}_\theta\big[\hat{g} - \mathbb{E}_\theta[\hat{g}]\big] = 0.
\]
Ainsi :
\[
\mathbb{E}_\theta\big[(\hat{g} - g(\theta))^2\big]
= \mathbb{E}_\theta\big[(\hat{g} - \mathbb{E}_\theta[\hat{g}])^2\big]
+ \big(\mathbb{E}_\theta[\hat{g}] - g(\theta)\big)^2,
\]
ce qui donne bien $R(\hat{g}, g(\theta)) = \mathrm{Var}_\theta(\hat{g}) + B(\hat{g}, g(\theta))^2$.
\end{proof}

\begin{remarque}
Si $\hat{g}$ est sans biais, alors $B(\hat{g}, g(\theta)) = 0$ et
\[
R(\hat{g}, g(\theta)) = \mathrm{Var}_\theta(\hat{g}).
\]
\end{remarque}

\paragraph{Interprétation :}
\begin{itemize}
  \item Le \textbf{biais} mesure l'erreur systématique faite par l'estimateur (erreur d'approximation).
  \item La \textbf{variance} mesure les fluctuations aléatoires de l'estimateur autour de sa moyenne (erreur d'estimation).
\end{itemize}

\paragraph{Remarques sur les estimateurs sans biais}

Si $\hat{g}$ est sans biais :
\[
R(\hat{g}, g(\theta)) = \mathrm{Var}_\theta(\hat{g}).
\]

On s'intéresse alors aux estimateurs sans biais ayant la variance minimale.  
On appelle un tel estimateur un \textbf{estimateur UVMB} (Uniformément de Variance Minimale parmi les estimateurs Sans Biais).  

Il existe des méthodes (exhaustivité, complétude, théorème de Lehmann-Scheffé) pour caractériser directement de tels estimateurs.

\begin{remarque}
Un estimateur sans biais n'est pas toujours le meilleur choix :  
un estimateur \textbf{avec biais} peut avoir un risque quadratique plus petit.  
De plus, il n'existe pas toujours d'estimateur sans biais pour une quantité donnée.
\end{remarque}

\begin{exemple}[Impossibilité d'un estimateur sans biais]
Soit $X \sim \mathcal{E}(\theta)$ une loi exponentielle de paramètre $\theta > 0$.  
Supposons qu'il existe un estimateur sans biais $\hat{\theta} = h(X)$ de $\theta$.  
Alors, pour tout $\theta > 0$, il faudrait :
\[
\int_0^\infty h(x)\, \theta e^{-\theta x}\, dx = \theta.
\]
Or, cette condition n'admet pas de solution mesurable $h$, donc un tel estimateur n'existe pas.
\end{exemple}

\begin{remarque}
Même si $\hat{\theta}$ est un estimateur sans biais de $\theta$, en général $g(\hat{\theta})$ n'est \emph{pas} un estimateur sans biais de $g(\theta)$ (sauf si $g$ est affine).  
Ainsi, la propriété « sans biais » n'est pas stable par transformation non linéaire.
\end{remarque}





\newpage

\section{Propriétés asymptotiques}

On se place dans le cadre de $n$ observations.  
On note $\hat{g}_n$ un estimateur de $g(\theta)$.

\smallskip

\begin{mdframed}
\begin{definition}[Consistance]
On dit que la suite $(\hat{g}_n)_{n \in \mathbb{N}}$ est \textbf{consistante} si, pour tout $\theta \in \Theta$,
\[
\hat{g}_n \xrightarrow[\mathbb{P}_\theta]{n \to +\infty} g(\theta).
\]
Elle est dite \textbf{fortement consistante} si, pour tout $\theta \in \Theta$,
\[
\hat{g}_n \xrightarrow[\text{p.s.}]{n \to +\infty} g(\theta).
\]
\end{definition}
\end{mdframed}

\begin{remarque}
\begin{itemize}
    \item La convergence doit avoir lieu pour \emph{tout} $\theta \in \Theta$.
    \item La consistance est une propriété « faible ». S'il existe un estimateur consistant, il en existe une infinité : si $\hat{g}_n$ est consistant et $(a_n)$ vérifie $a_n \to 1$, alors $a_n \hat{g}_n$ est aussi consistant.
    \item Les estimateurs non consistants sont à exclure : ils ne convergent pas vers la vraie valeur, quelle que soit la taille de l'échantillon.
\end{itemize}
\end{remarque}

\begin{exemple}[Consistance de la moyenne empirique]
Soit $(X_1, \ldots, X_n)$ un $n$-échantillon i.i.d.\ de loi $P_\theta$, avec $\mathbb{E}_\theta[X_1] = \mu(\theta) < +\infty$.  
La loi des grands nombres assure que
\[
\overline{X}_n = \frac{1}{n}\sum_{i=1}^n X_i \xrightarrow[\text{p.s.}]{n \to +\infty} \mu(\theta).
\]
Ainsi, $\overline{X}_n$ est un estimateur fortement consistant de $\mu(\theta)$.
\end{exemple}

\begin{definition}[Estimateur asymptotiquement sans biais]
On suppose que pour tout $\theta \in \Theta$ et tout $n \in \mathbb{N}$,
\[
\mathbb{E}_\theta\!\left[\|\hat{g}_n\|\right] < +\infty.
\]
On dit que $\hat{g}_n$ est \textbf{asymptotiquement sans biais} si, pour tout $\theta \in \Theta$,
\[
\mathbb{E}_\theta[\hat{g}_n] \xrightarrow{n \to +\infty} g(\theta).
\]
\end{definition}

\begin{remarque}
Asymptotiquement sans biais n'implique pas consistant, et réciproquement. En revanche, un estimateur sans biais et consistant est bien asymptotiquement sans biais.
\end{remarque}

\smallskip

\begin{mdframed}
\begin{definition}[Normalité asymptotique]
On dit que $\hat{g}_n$ est \textbf{asymptotiquement normal} si, pour tout $\theta \in \Theta$, il existe $\Sigma(\theta) > 0$ tel que
\[
\sqrt{n}\left(\hat{g}_n - g(\theta)\right) \xrightarrow[\text{loi}]{n \to +\infty} \mathcal{N}(0, \Sigma(\theta)).
\]
La quantité $\Sigma(\theta)$ est appelée la \textbf{variance asymptotique} de $\hat{g}_n$.
\end{definition}
\end{mdframed}

\begin{remarque}
\begin{itemize}
    \item Si $\hat{g}_n$ est asymptotiquement normal, alors $\hat{g}_n$ est consistant.
    \item La normalité asymptotique est fondamentale en pratique : elle permet de construire des intervalles de confiance et des tests asymptotiques.
    \item Plus généralement, on peut remplacer $\sqrt{n}$ par $a_n$ avec $a_n \to +\infty$. On dit alors que $1/a_n$ est la \textbf{vitesse de convergence} asymptotique.
\end{itemize}
\end{remarque}

\smallskip

\begin{mdframed}
\begin{proposition}
Soit $\hat{g}_n$ un estimateur asymptotiquement normal de $g(\theta)$ de variance asymptotique $\Sigma(\theta)$.
Soit $\hat{\Sigma}_n$ un estimateur consistant de $\Sigma(\theta)$, c'est-à-dire :
\[
\hat{\Sigma}_n \xrightarrow[\mathbb{P}_\theta]{n \to +\infty} \Sigma(\theta) \quad \forall\, \theta \in \Theta.
\]
Alors
\[
\frac{\sqrt{n}\left(\hat{g}_n - g(\theta)\right)}{\sqrt{\hat{\Sigma}_n}} \xrightarrow[\text{loi}]{n \to +\infty} \mathcal{N}(0, 1).
\]
\end{proposition}
\end{mdframed}

\begin{remarque}
Ce résultat est particulièrement utile en pratique : on remplace la variance asymptotique inconnue $\Sigma(\theta)$ par un estimateur $\hat{\Sigma}_n$, et la convergence en loi est préservée. Cela permet de construire des intervalles de confiance asymptotiques sans connaître $\theta$.
\end{remarque}

\begin{proof}
On écrit
\[
\frac{\sqrt{n}\left(\hat{g}_n - g(\theta)\right)}{\sqrt{\hat{\Sigma}_n}}
= \frac{\sqrt{n}\left(\hat{g}_n - g(\theta)\right)}{\sqrt{\Sigma(\theta)}}
\cdot \frac{\sqrt{\Sigma(\theta)}}{\sqrt{\hat{\Sigma}_n}},
\]
avec :
\begin{itemize}
    \item $\displaystyle\frac{\sqrt{n}\left(\hat{g}_n - g(\theta)\right)}{\sqrt{\Sigma(\theta)}} \xrightarrow[\text{loi}]{n \to +\infty} \mathcal{N}(0,1)$ par hypothèse,
    \item $\displaystyle\frac{\sqrt{\Sigma(\theta)}}{\sqrt{\hat{\Sigma}_n}} \xrightarrow[\mathbb{P}_\theta]{n \to +\infty} 1$ par consistance de $\hat{\Sigma}_n$.
\end{itemize}
Le \textbf{lemme de Slutsky} permet de conclure : le produit d'une suite convergeant en loi vers $Z$ et d'une suite convergeant en probabilité vers $1$ converge en loi vers $Z$.
\end{proof}

\section{Estimation de $\theta$ à partir d'un estimateur de $g(\theta)$}

\textbf{Question :} On dispose de $\hat{g}_n$, estimateur de $g(\theta)$. Comment en déduire un estimateur de $\theta$ ? Quelles propriétés possède-t-il ?

\smallskip

L'outil fondamental pour répondre à cette question est la \textbf{méthode delta}, que l'on énonce d'abord dans sa forme directe.

\begin{mdframed}
\begin{proposition}[Méthode delta]
Soit $\Theta$ un ouvert de $\mathbb{R}$. Soit $\hat{g}_n$ un estimateur de $g(\theta)$ tel que pour tout $\theta \in \Theta$,
\[
\sqrt{n}\left(\hat{g}_n - g(\theta)\right) \xrightarrow[\text{loi}]{n \to +\infty} \mathcal{N}(0, \sigma^2).
\]
Soit $\phi : g(\Theta) \to \mathbb{R}$ une fonction de classe $\mathcal{C}^1$. Alors, pour tout $\theta \in \Theta$,
\[
\sqrt{n}\left(\phi(\hat{g}_n) - \phi(g(\theta))\right) \xrightarrow[\text{loi}]{n \to +\infty} \mathcal{N}\!\left(0,\, \sigma^2 \left(\phi'(g(\theta))\right)^2\right).
\]
\end{proposition}
\end{mdframed}

\begin{remarque}
La méthode delta permet de propager la normalité asymptotique au travers d'une transformation régulière. Elle est très utilisée pour construire des estimateurs et quantifier leur précision.
\end{remarque}

\smallskip

On applique maintenant la méthode delta à la fonction $\phi = g^{-1}$, ce qui donne la version suivante.

\begin{mdframed}
\begin{proposition}[Méthode delta inverse]
Soit $\Theta$ un intervalle ouvert de $\mathbb{R}$.
\begin{itemize}
    \item Soit $\hat{g}_n$ estimateur de $g(\theta)$ tel que pour tout $\theta \in \Theta$,
    \[
    \sqrt{n}\left(\hat{g}_n - g(\theta)\right) \xrightarrow[\text{loi}]{n \to +\infty} \mathcal{N}(0, \sigma^2).
    \]
    \item On suppose que $g : \Theta \to g(\Theta)$ est un $\mathcal{C}^1$-difféomorphisme (bijective, $\mathcal{C}^1$, et $g^{-1} \in \mathcal{C}^1$). On définit :
    \[
    \hat{\theta}_n = g^{-1}(\hat{g}_n).
    \]
\end{itemize}
Alors :
\begin{itemize}
    \item $\hat{\theta}_n$ est consistant pour l'estimation de $\theta$.
    \item Pour tout $\theta \in \Theta$,
    \[
    \sqrt{n}\left(\hat{\theta}_n - \theta\right) \xrightarrow[\text{loi}]{n \to +\infty} \mathcal{N}\!\left(0,\, \frac{\sigma^2}{(g'(\theta))^2}\right).
    \]
\end{itemize}
\end{proposition}
\end{mdframed}

\begin{proof}
\begin{itemize}
    \item \textbf{Bonne définition de $\hat{\theta}_n$ :} $g$ est continue et bijective de $\Theta$ sur $g(\Theta)$, qui est un intervalle ouvert. Comme $\hat{g}_n \xrightarrow{\mathbb{P}_\theta} g(\theta) \in g(\Theta)$, on a
    \[
    \mathbb{P}\!\left(\hat{g}_n \in g(\Theta)\right) \xrightarrow{n \to +\infty} 1.
    \]
    Donc $\hat{\theta}_n = g^{-1}(\hat{g}_n)$ est bien défini avec une probabilité tendant vers $1$.

    \item \textbf{Consistance de $\hat{\theta}_n$ :} $g^{-1}$ est continue et $\hat{g}_n \xrightarrow{\mathbb{P}_\theta} g(\theta)$, donc par composition,
    \[
    \hat{\theta}_n = g^{-1}(\hat{g}_n) \xrightarrow{\mathbb{P}_\theta} g^{-1}(g(\theta)) = \theta.
    \]

    \item \textbf{Normalité asymptotique :} on applique la méthode delta à $\phi = g^{-1}$. On rappelle que
    \[
    (g^{-1})'(y) = \frac{1}{g'(g^{-1}(y))},
    \]
    d'où $(g^{-1})'(g(\theta)) = \dfrac{1}{g'(\theta)}$.

    Par la méthode delta :
    \[
    \sqrt{n}\left(\hat{\theta}_n - \theta\right) = \sqrt{n}\left(g^{-1}(\hat{g}_n) - g^{-1}(g(\theta))\right) \xrightarrow[\text{loi}]{n \to +\infty} \mathcal{N}\!\left(0,\, \frac{\sigma^2}{(g'(\theta))^2}\right).
    \]
\end{itemize}
\end{proof}
 %FINIS

\chapter*{TD 1 - Révisions de probabilité}

\setcounter{exercice}{0}

\begin{mdframed}
\begin{exercice}
\begin{enumerate}
    \item Soit $X \sim \mathcal{E}(\lambda)$. Donner la loi de la variable $Y = 1 + \lfloor X \rfloor$, où $\lfloor \cdot \rfloor$ désigne la partie entière.
    \item Soit $U \sim \mathcal{U}([0, 1])$. Donner la loi de $Y = U^2$.
    \item Soit $U \sim \mathcal{U}([0, 1])$. Donner la loi de $Y = \lfloor 1/U \rfloor$.
    \item Soit $X \sim \mathcal{N}(0, 1)$. Donner la loi de $Y = X^2$.
\end{enumerate}
\end{exercice}
\end{mdframed}

\smallskip

\begin{mdframed}
\begin{exercice}
Soit $X_1, \dots, X_n$ des v.a. indépendantes, centrées, telles que pour tout $i \in \{1, \dots, n\}$, il existe des réels $c_i > 0$ vérifiant, $|X_i| \le c_i$ p.s.. On note
\begin{equation*}
S_n = \sum_{i=1}^n X_i, \quad \text{et} \quad C_n = \sum_{i=1}^n c_i^2.
\end{equation*}
L'objectif est de prouver l'inégalité suivante (appelée inégalité de Hoeffding) : pour tout $x \ge 0$,
\begin{equation*}
\mathbb{P}(|S_n| \ge x) \le 2 \exp \left( -\frac{x^2}{2C_n} \right).
\end{equation*}
\begin{enumerate}
    \item Soit $X$ une v.a.r. On définit pour tout $t \ge 0$ et tout $x \ge 0$,
    \begin{equation*}
    \ell_X(t) = \log \mathbb{E}\left[e^{tX}\right], \quad \ell_X^*(x) = \sup_{t \ge 0} \{xt - \ell_X(t)\}.
    \end{equation*}
    Montrer que pour tout $x \ge 0$,
    \begin{equation*}
    \mathbb{P}(X \ge x) \le \exp\left(-\ell_X^*(x)\right).
    \end{equation*}
    
    \item Montrer que pour tout $t \ge 0$,
    \begin{equation*}
    \ell_{S_n}(t) = \sum_{i=1}^n \ell_{X_i}(t).
    \end{equation*}
    
    \item Soit $X$ une v.a. centrée telle que $|X| \le 1$ p.s. Soit $\theta$ une v.a. de loi de Rademacher (i.e. prenant les valeurs $-1$ et $1$ avec probabilité $1/2$). Vérifier que pour toute fonction convexe $\phi : \mathbb{R} \to \mathbb{R}$,
    \begin{equation*}
    \mathbb{E}[\phi(X)] \le \mathbb{E}[\phi(\theta)].
    \end{equation*}
    En déduire que pour tout $x \ge 0$,
    \begin{equation*}
    \ell_X(t) \le \frac{t^2}{2}.
    \end{equation*}
    
    \item Montrer que pour tout $x \ge 0$,
    \begin{equation*}
    \mathbb{P}(S_n \ge x) \le \exp \left(-\frac{x^2}{2C_n}\right).
    \end{equation*}
    \item Conclure.
\end{enumerate}
\end{exercice}
\end{mdframed}

\smallskip

\begin{mdframed}
\begin{exercice}
\begin{enumerate}
    \item Soit $X$ une v.a. à valeurs dans $\mathbb{N}$. Montrer que
    \begin{equation*}
    \mathbb{E}[X] = \sum_{n \ge 0} \mathbb{P}(X > n).
    \end{equation*}
    \item Soit $X$ une v.a. positive p.s. Montrer que pour tout $\alpha > 0$,
    \begin{equation*}
    \mathbb{E}[X^\alpha] = \alpha \int_0^{+\infty} t^{\alpha-1} \mathbb{P}(X > t) dt.
    \end{equation*}
\end{enumerate}
\end{exercice}
\end{mdframed}

\smallskip

\begin{mdframed}
\begin{exercice}
Soit $(X_n)_n$ une suite de v.a.r. i.i.d. de carrés intégrables. Démontrer que
\begin{equation*}
\mathbb{E}[\bar{X}_n] = \mathbb{E}[X_1] \quad \text{et} \quad \text{Var}(\bar{X}_n) = \frac{\text{Var}(X_1)}{n}.
\end{equation*}
\end{exercice}
\end{mdframed}

\smallskip

\begin{mdframed}
\begin{exercice}
Soit $(U_i)_{i \in \mathbb{N}}$ une suite de variables aléatoires indépendantes et identiquement distribuées de loi uniforme sur $[a, b]$, $a < b$.
\begin{enumerate}
    \item Déterminer la loi de $M_n = \sup_{1 \le i \le n} U_i$.
    \item Même question pour l'infimum : $m_n = \inf_{1 \le i \le n} U_i$.
    \item Étudier $\mathbb{P}(n(b - M_n) > x)$ lorsque $n$ tend vers l'infini. Que peut-on en déduire ?
\end{enumerate}
\end{exercice}
\end{mdframed}

\smallskip

\begin{mdframed}
\begin{exercice}
Soit $(X_i)_{i \ge 1}$ une suite de variables aléatoires indépendantes, identiquement distribuées, centrées et réduites. Par le théorème central limite on a
\begin{equation*}
Z_n = \frac{1}{\sqrt{n}} \sum_{i=1}^n X_i \xrightarrow[n \to +\infty]{\text{loi}} Z \sim \mathcal{N}(0, 1).
\end{equation*}
Le but de cet exercice est de montrer que $Z_n$ ne peut pas converger en probabilité.
\begin{enumerate}
    \item Décomposer la variable $Z_{2n}$ en fonction de $Z_n$ et d'une variable indépendante de $Z_n$.
    \item Calculer la fonction caractéristique de $Z_{2n} - Z_n$ et montrer que cette différence converge en loi.
    \item En raisonnant par l'absurde, en déduire que $Z_n$ ne converge pas en probabilité. \textit{Indication : on pourra démontrer le résultat suivant : }
    \begin{quote}
    \textbf{Lemme 1} Soit $(U_n)_n$ une suite de v.a.r. convergeant en probabilité vers $U$ et $(V_n)_n$ une suite de v.a.r. convergeant en probabilité vers $V$. Soit $a \in \mathbb{R}$. Alors
    \begin{equation*}
    aU_n + V_n \xrightarrow[n \to +\infty]{\mathbb{P}} aU + V.
    \end{equation*}
    \end{quote}
\end{enumerate}
\end{exercice}
\end{mdframed}

\smallskip

\begin{mdframed}
\begin{exercice}
On rappelle que l'on dit qu'un couple de v.a.r. $Z = (X, Y)$ suit la loi uniforme sur une partie (mesurable) $B \subset \mathbb{R}^2$ si la densité de $Z$ est donnée par $f_Z(z) = \frac{\mathbf{1}_B(z)}{\text{Vol}(B)}$, où $\text{Vol}(B)$ est le volume de $B$.
\begin{enumerate}
    \item Soit $X, Y$ un point tiré au hasard (uniformément) sur le carré unité de $\mathbb{R}^2$. Montrer que $X$ et $Y$ sont indépendantes et de loi uniforme sur $[0, 1]$.
    \item Soit $(X, Y)$ un point tiré au hasard (uniformément) sur le disque unité de $\mathbb{R}^2$. Déterminer la loi marginale de $X$.
\end{enumerate}
\end{exercice}
\end{mdframed}

\smallskip

\begin{mdframed}
\begin{exercice}
\begin{enumerate}
    \item On coupe un bâton au hasard en trois morceaux en utilisant deux v.a. indépendantes et uniformes sur $[0, 1]$ pour déterminer les points de coupe. Vérifier que les longueurs des trois morceaux ainsi obtenus sont des v.a. de même loi. Sont-elles indépendantes ? Quelle est la probabilité de pouvoir fabriquer un triangle avec ces trois morceaux ?
    \item Deux amis se donnent rendez-vous entre 12h et 13h, et arrivent indépendamment uniformément entre ces deux horaires. Calculer le temps moyen d'attente du premier arrivé.
\end{enumerate}
\end{exercice}
\end{mdframed}

\smallskip

\begin{mdframed}
\begin{exercice}
Écrire un programme pour simuler $n$ v.a. i.i.d., $X_1, \dots, X_n$, telles que :
\begin{enumerate}
    \item $X_1$ est une v.a. discrète à valeurs dans $\{w_1, \dots, w_r\}$ avec pour tout $1 \le i \le r$,
    \begin{equation*}
    \mathbb{P}(X_1 = w_i) = p_i.
    \end{equation*}
    Les valeurs $w_1, \dots, w_r$, et $p_1, \dots, p_r$ sont des arguments du programme.
    \item $X_1$ suit une loi de Laplace, de densité sur $\mathbb{R}$ donnée par $x \mapsto \frac{1}{2} e^{-|x|}$.
    \item $X_1$ suit une loi de Cauchy, de densité sur $\mathbb{R}$ donnée par $x \mapsto \frac{1}{\pi} \frac{1}{1+x^2}$.
\end{enumerate}
\end{exercice}
\end{mdframed}

\smallskip

\begin{mdframed}
\begin{exercice}
Soit $Z$ une variable aléatoire gaussienne centrée réduite.
\begin{enumerate}
    \item Montrer que pour tout $\lambda \in \mathbb{R}$, $M_Z(\lambda) := \mathbb{E}[e^{\lambda Z}] = e^{\lambda^2/2}$.
    \item En déduire que $\mathbb{E}[Z^n] = 0$ si $n$ est impair et $\mathbb{E}[Z^{2k}] = \frac{(2k)!}{2^k k!}$ pour tout $k \ge 0$.
    \item Conclure que $\varphi_Z(t) := \mathbb{E}[e^{itZ}] = e^{-t^2/2}$, pour tout $t \in \mathbb{R}$.
\end{enumerate}
\end{exercice}
\end{mdframed}

\smallskip

\begin{mdframed}
\begin{exercice}
Soit $X_n, n \in \mathbb{N}$, et $X$ des variables aléatoires à valeurs dans $\mathbb{N}$.
\begin{enumerate}
    \item On suppose maintenant que pour tout $n \in \mathbb{N}$, $X_n \sim \mathcal{B}(n, p_n)$, avec $p_n \in [0, 1]$ tel que $np_n \xrightarrow[n \to +\infty]{} \lambda > 0$. Montrer que $\mathbb{P}(X_n = k) \xrightarrow[n \to +\infty]{} \mathbb{P}(X = k)$ avec $X$ de loi de Poisson de paramètre $\lambda$. Conclure.
    \item Retrouver le résultat précédent en utilisant le théorème de Lévy.
\end{enumerate}
\end{exercice}
\end{mdframed}

\smallskip

\begin{mdframed}
\begin{exercice}
Dans cet exercice, on présente une preuve rapide de la loi forte des grands nombres sous une hypothèse de moment d'ordre 4 fini.\\
Soit $(X_n)_{n \ge 1}$ une suite de v.a.r. i.i.d. telles que $\mathbb{E}[|X_1|^4] < \infty$. On pose pour tout $n \in \mathbb{N}$, $S_n = \sum_{k=1}^n X_k$ et on note $m = \mathbb{E}[X_1]$.

\begin{enumerate}
    \item Montrer qu'il existe une constante $K > 0$ telle que, pour tout $n \ge 1$
    \begin{equation*}
    \mathbb{E}[(S_n - nm)^4] \le Kn^2.
    \end{equation*}
    \item En déduire que $\frac{S_n}{n} \xrightarrow[n \to +\infty]{\text{p.s.}} m$.
\end{enumerate}
\end{exercice}
\end{mdframed}


\chapter{Méthode de estimation}

\section{Méthode des moments}

\smallskip

\begin{mdframed}
\begin{definition}[Estimateur par la méthode des moments]

\begin{itemize}
    \item Soit un échantillon $X_1, \ldots, X_n$ de loi $\mathbb{P}_\theta$, $\theta \in \Theta$.
    \item Soit $\phi$ une fonction telle que $\mathbb{E}_\theta\left[\|\phi(X_1)\|\right] < +\infty$.
\end{itemize}
On souhaite estimer
\[
g(\theta) = \mathbb{E}_\theta[\phi(X_1)].
\]
On appelle \textbf{estimateur par la méthode des moments} :
\[
\hat{g}_n = \frac{1}{n} \sum_{i=1}^n \phi(X_i).
\]
\end{definition}
\end{mdframed}


\begin{notation}
\begin{itemize}
    \item $\hat{g}_n$ est consistant (fortement) par la loi des grands nombres.
    \item $\hat{g}_n$ est sans biais (par linéarité de l'espérance).
    \item Plus généralement, si $\hat{g}_n$ est construit par la méthode des moments comme ci-dessus, et si $h$ est une fonction, on dira aussi que
    \[
    h(\hat{g}_n)
    \]
    est un estimateur par la \textbf{méthode des moments} pour l'estimation de $h(g(\theta))$.
\end{itemize}
\end{notation}


\begin{remarque}
À priori, on ne connaît pas les propriétés de cet estimateur (biais, variance, normalité asymptotique). Il faudra étudier ces propriétés au cas par cas, éventuellement en utilisant la méthode delta.
\end{remarque}

\begin{exemple}
\begin{enumerate}
    \item \textbf{Probabilité d'un événement}
    
    Soit $A$ un événement. On souhaite estimer
    \[
    g(\theta) = \mathbb{P}_\theta(X_1 \in A) = \mathbb{E}_\theta[\mathbb{1}_{X_1 \in A}],
    \]
    où
    \[
    \mathbb{1}_{X_1 \in A} = \begin{cases}
    1 & \text{si } X_1 \in A, \\
    0 & \text{sinon}.
    \end{cases}
    \]
    
    Par la méthode des moments :
    \[
    \hat{g}_{A,n} = \frac{1}{n} \sum_{i=1}^n \mathbb{1}_{X_i \in A}.
    \]
    
    Cet estimateur est la fréquence empirique de l'événement $A$.
    
    \item \textbf{Loi uniforme}
    
    Soit $X_i$ pour $i \in [ 1, n ]$ de loi $\mathcal{U}([\theta-1 , \theta+1])$ avec $\theta \in \mathbb{R}$.
    
    Alors
    \[
    \mathbb{E}_\theta[X_1] = \theta = g(\theta).
    \]
    
    Par la méthode des moments :
    \[
    \hat{g}_n = \overline{X}_n = \frac{1}{n} \sum_{i=1}^n X_i,
    \]
\end{enumerate}
\end{exemple}




\section{Méthode du maximum de vraisemblance}

On se place dans le cadre de l'estimation de $g(\theta)=\theta$.

\smallskip

\begin{mdframed}
\begin{definition}
On appelle vraisemblance associée à l'observation $X$ la fonction

$$
L : 
\begin{cases}
\Theta \to \mathbb{R} \\
\theta \mapsto L(\theta,X)
\end{cases}
$$

définie par

\begin{itemize}
\item Si $X$ v.a. discrète

$$
L(\theta,X) = P_\theta (X) = \mathbb{P}_\theta (X=x)
$$
\item Si $X$ v.a. à densité

$$
L(\theta,X) = f_\theta(X)
$$

où $f_\theta$ est la densité de $X$ sous $\mathbb{P}_\theta$
\end{itemize}
\end{definition}
\end{mdframed}

\begin{remarque}
Si $X=(X_1,\cdots,X_n)$ est un $n$-échantillon i.i.d.

\begin{itemize}
\item $X_i$ discret :

$$
L(\theta,X) = \prod_{i=1}^{n} P_\theta (X_i)
$$
\item $X_i$ à densité :

$$
L(\theta,X) = \prod_{i=1}^{n} f_\theta(X_i)
$$
\end{itemize}
\end{remarque}

\begin{exemple}
On considère $X \sim \mathcal{B}(15,p)$ avec $p$ inconnue. On observe $x=5$.

La fonction de vraisemblance est

$$
L(p,5) = \mathbb{P}_p (X=5) = \binom{15}{5} p^5 (1-p)^{10}
$$

C'est la probabilité d'observer $5$ quand le paramètre est $p$.

\begin{center}
\begin{tabular}{|c|c|c|c|c|c|}
\hline
$p$ & 0,1 & 0,2 & 0,3 & 0,4 & 0,5 \\
\hline
$L(p,5)$ & 0,01 & 0,1 & 0,21 & 0,19 & 0,09 \\
\hline
\end{tabular}
\end{center}

Quand $p=0,3$ il y a $21$ chances sur $100$ d'observer $x=5$.

Il est donc plus « vraisemblable » que le vrai $p$ soit égal à $0,3$ plutôt qu'à $0,1$.

En suivant ce raisonnement, la valeur de $p$ la plus vraisemblable est celle qui maximise la probabilité d'observer $5$.

C'est donc la valeur qui maximise la vraisemblance.

Pour trouver ce maximum, on utilise la log-vraisemblance :

\begin{align*}
l(p,x) &= \ln L(p,x) = \ln \binom{15}{x} + x\ln(p) + (15-x)\ln(1-p)
\end{align*}

On dérive par rapport à $p$ :

\begin{align*}
\frac{\partial l}{\partial p}(p,x)
&= \frac{x}{p} - \frac{15-x}{1-p}\\
&= \frac{x(1-p) - (15-x)p}{p(1-p)}\\
&= \frac{x-15p}{p(1-p)}
\end{align*}

Cette dérivée s'annule en $p=\frac{x}{15}$.

On vérifie qu'il s'agit bien d'un maximum :

$$
\frac{\partial^2 l}{\partial p^2}(p,x) = -\frac{x}{p^2} - \frac{15-x}{(1-p)^2} < 0
$$

Ici avec $x=5$ on a $\hat{p}=\frac{5}{15}=\frac{1}{3}$.

La valeur de $p$ la plus vraisemblable au vu de l'observation $x=5$ est $p=\frac{1}{3}$.
\end{exemple}

\smallskip

\begin{mdframed}
\begin{definition}
On appelle estimateur du maximum de vraisemblance (E.M.V.) de $\theta$ toute quantité $\hat{\theta}$ satisfaisant :

$$
\hat{\theta}^{\text{EMV}} \in \arg \max_{\theta \in \Theta} L(\theta,X)
$$

Souvent on considère la log-vraisemblance :

$$
\hat{\theta} \in \arg \max_{\theta \in \Theta} l(\theta,X)
$$

où $l(\theta,X) = \ln L(\theta,X)$.
\end{definition}
\end{mdframed}

\begin{proposition}[Méthode pratique de calcul]
Si $L(\theta, X)$ (ou $l(\theta,X)$) est dérivable par rapport à $\theta$ et si le maximum est atteint à l'intérieur de $\Theta$, alors $\hat{\theta}$ vérifie l'équation de vraisemblance :
$$
\frac{\partial l}{\partial \theta}(\hat{\theta}, X) = 0
$$
avec la condition du second ordre :
$$
\frac{\partial^2 l}{\partial \theta^2}(\hat{\theta}, X) < 0
$$
ou, en dimension supérieure, la hessienne doit être définie négative.
\end{proposition}

\begin{theoreme}[Principe d'invariance]
Si $\hat{\theta}$ est un EMV de $\theta$ et $g : \Theta \to \mathbb{R}$ une fonction, alors $g(\hat{\theta})$ est un EMV de $g(\theta)$.
\end{theoreme}

\begin{exemple}
Soit $X_1, \ldots, X_n$ un $n$-échantillon i.i.d. de loi $\mathcal{N}(\theta, 1)$ avec $\theta \in \mathbb{R}$ inconnu.
La vraisemblance est :
$$
L(\theta, X) = \prod_{i=1}^{n} \frac{1}{\sqrt{2\pi}} e^{-\frac{(X_i-\theta)^2}{2}} = (2\pi)^{-n/2} e^{-\frac{1}{2}\sum_{i=1}^{n}(X_i-\theta)^2}
$$
La log-vraisemblance est :
$$
l(\theta, X) = -\frac{n}{2}\ln(2\pi) - \frac{1}{2}\sum_{i=1}^{n}(X_i-\theta)^2
$$
On dérive :
$$
\frac{\partial l}{\partial \theta}(\theta, X) = \sum_{i=1}^{n}(X_i-\theta) = n\bar{X}_n - n\theta
$$
Cette dérivée s'annule en $\theta = \bar{X}_n$. Comme $\frac{\partial^2 l}{\partial \theta^2} = -n < 0$, il s'agit bien d'un maximum.

Donc l'EMV de $\theta$ est : $\hat{\theta} = \bar{X}_n$.

Par le principe d'invariance, l'EMV de $\theta^2$ est $(\bar{X}_n)^2$.
\end{exemple}

\begin{exemple}
Soit $X_1, \ldots, X_n$ un $n$-échantillon i.i.d. de loi exponentielle $\mathcal{E}(\lambda)$ de densité $f_\lambda(x) = \lambda e^{-\lambda x} \mathbb{1}_{x \geq 0}$, avec $\lambda > 0$ inconnu.

La vraisemblance est :
$$
L(\lambda, X) = \prod_{i=1}^{n} \lambda e^{-\lambda X_i} = \lambda^n e^{-\lambda \sum_{i=1}^{n} X_i}
$$
La log-vraisemblance est :
$$
l(\lambda, X) = n\ln(\lambda) - \lambda \sum_{i=1}^{n} X_i
$$
On dérive :
$$
\frac{\partial l}{\partial \lambda}(\lambda, X) = \frac{n}{\lambda} - \sum_{i=1}^{n} X_i
$$
Cette dérivée s'annule en $\lambda = \frac{n}{\sum_{i=1}^{n} X_i} = \frac{1}{\bar{X}_n}$.

Comme $\frac{\partial^2 l}{\partial \lambda^2} = -\frac{n}{\lambda^2} < 0$, il s'agit bien d'un maximum.

Donc l'EMV de $\lambda$ est : $\hat{\lambda} = \frac{1}{\bar{X}_n}$.
\end{exemple}

\begin{remarque}[Avantages et inconvénients]
\textbf{Avantages :}
\begin{itemize}
\item Méthode générale applicable à de nombreux modèles statistiques.
\item Interprétation intuitive : on choisit le paramètre qui rend l'observation la plus probable.
\item Bonnes propriétés asymptotiques (convergence, normalité asymptotique, efficacité).
\item Le principe d'invariance facilite l'estimation de fonctions de paramètres.
\end{itemize}

\textbf{Inconvénients :}
\begin{itemize}
\item L'EMV peut ne pas être unique.
\item L'EMV peut ne pas exister.
\item L'EMV peut être difficile à exhiber (pas de forme explicite).
\item Peut être biaisé en échantillon fini (même si asymptotiquement sans biais).
\end{itemize}
\end{remarque} %FINIS

\chapter*{TD 2 - Méthodes d'estimation - 1ère partie}

\setcounter{exercice}{0}

\begin{mdframed}
\begin{exercice}
Les familles de lois suivantes sont-elles identifiables ?
\begin{enumerate}
    \item $\{\mathcal{E}(\lambda), \, \lambda \in (0, +\infty)\}$.
    \item $\{\mathcal{E}(\lambda_1 \lambda_2), \, (\lambda_1, \lambda_2) \in (0, +\infty)^2\}$.
    \item $\{\mathcal{B}(p), \, p \in (0, 1)\}$.
    \item $\{\mathcal{N}(0, \sigma^2), \, \sigma \in \mathbb{R}\}$.
    \item $\{\mathcal{N}(0, \sigma^2), \, \sigma^2 \in (0, +\infty)\}$.
\end{enumerate}
\end{exercice}
\end{mdframed}

\smallskip

\begin{mdframed}
\begin{exercice}
Soit $(X_i)_{i \ge 1}$ une suite de variables aléatoires indépendantes et identiquement distribuées. On note $m = \mathbb{E}[X_1]$ et $\sigma^2 = \text{Var}(X_1)$. Vérifier que
\begin{equation*}
\bar{X}_n \quad \text{et} \quad S_n = \frac{1}{n} \sum_{i=1}^n (X_i - \bar{X}_n)^2
\end{equation*}
sont des estimateurs consistants de $m$ et de $\sigma^2$. Calculer leur biais et en déduire un estimateur sans biais de $\sigma^2$.
\end{exercice}
\end{mdframed}

\smallskip

\begin{mdframed}
\begin{exercice}
Soit $(X_i)_{i \ge 1}$ une suite de variables aléatoires indépendantes et identiquement distribuées. On note $m = \mathbb{E}[X_1]$ et $\sigma^2 = \text{Var}(X_1)$. On propose comme estimateurs de $m$, $\bar{X}_n$ et
\begin{equation*}
T_{n,1} = \frac{1}{2}(X_n + X_{n-1}).
\end{equation*}

\begin{enumerate}
    \item Montrer que $T_{n,1}$ est sans biais.
    \item Lequel des deux estimateurs $\bar{X}_n$ et $T_{n,1}$ choisiriez-vous pour estimer $m$ ?
    \item Que pensez-vous de l'estimateur $T_{n,2} = 0$ pour l'estimation de $m$ ?
\end{enumerate}
\end{exercice}
\end{mdframed}

\smallskip

\begin{mdframed}
\begin{exercice}
Lors du contrôle d’une chaîne de médicaments, on s’intéresse au nombre de comprimés défectueux dans un lot. On note $X_i$ le nombre de comprimés défectueux dans le lot $i$. On suppose que ces variables aléatoires sont i.i.d. de loi de probabilité inconnue, d’espérance $\mu$ et de variance $\sigma^2$.

Le contrôle effectué sur 20 lots choisis au hasard ont donné les résultats suivants :

$$
1\quad 0 \quad 0 \quad 3 \quad 2 \quad 0 \quad 5 \quad 2 \quad 0 \quad 0 \quad 1 \quad 2 \quad 1 \quad 3 \quad 0 \quad 1 \quad 0 \quad 0 \quad 2 \quad 7
$$

\begin{enumerate}
    \item On considère quatre estimateurs pour $\mu$ :
    $$
    T_1 = X_1 \quad \frac{X_1 + X_2}{2} \quad T_3 = \frac{X_1 + X_2}{3} \quad T_4 = \frac{\sum_i^{20} X_i}{20}
    $$
    \begin{enumerate}
        \item[(a)] Calculer le biais et la variance de chacun des estimateurs.
        \item[(b)] Quel est le meilleur estimateur ?
        \item[(c)] Quelle est l’estimation correspondante (c’est-à-dire sa valeur observée) ?
    \end{enumerate}
    \item Proposer un estimateur de la probabilité qu’un lot contienne au moins un comprimé défectueux et calculer
l’estimation correspondante.
\end{enumerate}
\end{exercice}
\end{mdframed}

\smallskip

\begin{mdframed}
\begin{exercice}
Soit $X_1, \cdots , X_n$ un échantillon i.i.d. de loi de probabilité $\mathbb{P}(X = 1) = \theta = 1 - \mathbb{P}(X = -1)$. On se propose d’estimer le paramètre $\theta \in [0 , 1]$ par l’estimateur $\hat{\theta} = \frac{\bar{X} + 1}{2}$.

\begin{enumerate}
    \item Calculer $\mathbb{E}[X]$ et en déduire que $\hat{\theta}$ est un estimateur sans biais de $\theta$.
    \item Montrer que $\hat{\theta}$ est un estimateur consistant de $\theta$.
    \item Déterminer le risque quadratique de $\hat{\theta}$.
\end{enumerate}
\end{exercice}
\end{mdframed}

\smallskip

\begin{mdframed}
\begin{exercice}
Soit $(X_i)_{i \ge 1}$ une suite de variables aléatoires indépendantes et identiquement distribuées de loi normale de moyenne $m$ et de variance $\sigma^2$.

\begin{enumerate}
    \item Calculer la loi de $\bar{X}_n$.
    \item On suppose que $\sigma$ est connu.
    \begin{enumerate}
        \item[i--] Donner un majorant de $\mathbb{P}_m\left( |\bar{X}_n - m|/\sigma > t \right)$ au moyen de l'inégalité de Bienaymé-Tchebychev. En déduire un majorant du nombre d'observations minimal nécessaire pour que cette probabilité soit inférieure à $0.05$. Donner la valeur pour $t = 0.01$.
        \item[ii--] Calculer ce nombre minimal en utilisant la vraie loi. Donner l'application numérique associée. Que remarquez-vous ?
        \item[iii--] En déduire la valeur de $\mathbb{P}_m\left( m \in \left[ \bar{X}_n - 1.96 \frac{\sigma}{\sqrt{n}}, \, \bar{X}_n + 1.96 \frac{\sigma}{\sqrt{n}} \right] \right)$.
        \item[iv--] Comment s'interprète l'intervalle $I = \left[ \bar{X}_n - 1.96 \frac{\sigma}{\sqrt{n}}, \, \bar{X}_n + 1.96 \frac{\sigma}{\sqrt{n}} \right]$ ?
    \end{enumerate}
\end{enumerate}
\end{exercice}
\end{mdframed}

\smallskip

\begin{mdframed}
\begin{exercice}
Soit $(X_i)_{i \ge 1}$ une suite de variables aléatoires indépendantes et identiquement distribuées de loi exponentielle $\mathcal{E}(\theta)$, $\theta > 0$.

\begin{enumerate}
    \item Un statisticien propose d'estimer $\theta$ par :
    \begin{equation*}
    \hat{\theta}_n = \frac{1}{\bar{X}_n}.
    \end{equation*}
    \begin{enumerate}
        \item[i--] Justifier que $\hat{\theta}_n$ est bien définie.
        \item[ii--] Calculer $\mathbb{E}_\theta[X_1]$ et expliquer le choix du statisticien.
        \item[iii--] Donner la loi limite de $\sqrt{n}(\hat{\theta}_n - \theta)$.
        \item[iv--] Donner la loi de $\sum_{i=1}^n X_i$. Calculer $\mathbb{E}_\theta[\hat{\theta}_n]$, $\text{Var}_\theta(\hat{\theta}_n)$, et $\mathbb{E}_\theta[(\hat{\theta}_n - \theta)^2]$.
    \end{enumerate}
    
    \item Un statisticien propose d'utiliser
    \begin{equation*}
    \hat{\theta}_n^{(2)} = \frac{n-1}{n} \hat{\theta}_n.
    \end{equation*}
    \begin{enumerate}
        \item[i--] L'estimateur $\hat{\theta}_n^{(2)}$ vérifie-t-il toujours des propriétés asymptotiques similaires à $\hat{\theta}_n$ ?
        \item[ii--] Calculer $\mathbb{E}_\theta[\hat{\theta}_n^{(2)}]$, $\text{Var}_\theta(\hat{\theta}_n^{(2)})$, et $\mathbb{E}_\theta[(\hat{\theta}_n^{(2)} - \theta)^2]$. Quel estimateur préférez-vous ?
    \end{enumerate}
\end{enumerate}
\end{exercice}
\end{mdframed}

\smallskip

\begin{mdframed}
\begin{exercice}
Soit $(X_i)_{i \ge 1}$ une suite de variables aléatoires indépendantes et identiquement distribuées de loi de densité sur $\mathbb{R}$ donnée par
\begin{equation*}
f_\theta(x) = \frac{2x}{\theta^2} \mathbf{1}_{[0, \theta]}(x),
\end{equation*}
où $\theta > 0$ est un paramètre inconnu.

\begin{enumerate}
    \item On pose $X_{(n)} = \max_{1 \le i \le n} X_i$.
    \begin{enumerate}
        \item[i--] Donner la densité de $X_{(n)}$.
        \item[ii--] Calculer $\mathbb{E}_\theta[X_{(n)}]$, $\mathbb{E}_\theta[X_{(n)}^2]$, puis en déduire $\text{Var}_\theta(X_{(n)})$.
        \item[iii--] Montrer que $X_{(n)}$ est consistant pour l'estimation de $\theta$.
    \end{enumerate}
    \item Calculer $\mathbb{E}_\theta[X_1]$ puis en déduire un estimateur $\hat{\theta}_n$ consistant de $\theta$.
    \item Qui de $X_{(n)}$ ou $\hat{\theta}_n$ choisiriez-vous pour estimer $\theta$ ?
\end{enumerate}
\end{exercice}
\end{mdframed}

\smallskip

\begin{mdframed}
\begin{exercice}
Soit $(X_i)_{i \ge 1}$ une suite de variables aléatoires indépendantes et identiquement distribuées de loi de densité sur $\mathbb{R}$ donnée par
\begin{equation*}
f_\theta(x) = \frac{1}{\theta} \exp\left( -\frac{x - \theta}{\theta} \right) \mathbf{1}_{[\theta, +\infty[}(x),
\end{equation*}
où $\theta > 0$ est un paramètre inconnu.

\begin{enumerate}
    \item Montrer que $X_1/\theta - 1$ suit une loi exponentielle $\mathcal{E}(1)$.
    \item 
    \begin{enumerate}
        \item[i--] Déterminer la fonction de répartition de $X_{(1)} = \min_{1 \le i \le n} X_i$.
        \item[ii--] Calculer le risque quadratique de $X_{(1)}$ pour l'estimation de $\theta$ puis analyser la consistance de cet estimateur.
    \end{enumerate}
    \item 
    \begin{enumerate}
        \item[i--] Estimer $\theta$ par la méthode des moments.
        \item[ii--] Quelle est la loi limite de l'estimateur $\hat{\theta}$ ainsi obtenu ? Étudier également sa convergence en moyenne quadratique.
    \end{enumerate}
    \item Comparer les deux estimateurs.
\end{enumerate}
\end{exercice}
\end{mdframed}

\smallskip
\chapter*{TD 3 - Méthodes d'estimation - 2ème partie}

\setcounter{exercice}{0}

\begin{mdframed}
\begin{exercice}
Soit $(X_1, \dots, X_n)$ un $n$-échantillon d'une loi de Bernoulli de paramètre $\theta \in \, ]0, 1[$.
\begin{enumerate}
    \item Déterminer un estimateur de $\theta$ par la méthode des moments et donner ses propriétés.
    \item Déterminer l'estimateur du maximum de vraisemblance de $\theta$.
\end{enumerate}
\end{exercice}
\end{mdframed}

\smallskip

\begin{mdframed}
\begin{exercice}
Soit $X$ une v.a. de loi exponentielle $\mathcal{E}(\theta)$, $\theta > 0$. Soit $(X_1, \dots, X_n)$ un $n$-échantillon de même loi que $X$. On veut estimer $\theta$.
\begin{enumerate}
    \item Calculer l'estimateur du maximum de vraisemblance $\hat{\theta}_n$ de $\theta$.
    \item En reprenant l'exercice 5, rappeler le biais de $\hat{\theta}_n$, un estimateur sans biais $\hat{\theta}_n^*$ et son risque quadratique.
    \item Déterminer un estimateur $\hat{\theta}_n^\circ$ qui minimise le risque quadratique parmi les estimateurs s'écrivant :
    \begin{equation*}
    \hat{\theta}_n(c) = \frac{c}{\sum_{i=1}^n X_i},
    \end{equation*}
    où $c > 0$.
    \item Lequel préférez-vous ?
\end{enumerate}
\end{exercice}
\end{mdframed}

\smallskip

\begin{mdframed}
\begin{exercice}
Soit $X$ une v.a. de loi de densité donnée pour tout $x \in \mathbb{R}$, par
\begin{equation*}
f_\theta(x) = \exp(-(x - \theta))\mathbf{1}_{[\theta, +\infty[} (x),
\end{equation*}
où $\theta > 0$ est un paramètre inconnu. Soit $(X_1, \dots, X_n)$ un $n$-échantillon de même loi que $X$. On veut estimer $\theta$.
\begin{enumerate}
    \item Calculer l'estimateur du maximum de vraisemblance de $\theta$ que l'on notera $\hat{\theta}_n$.
    \item Calculer pour tout $t \in \mathbb{R}$, $\mathbb{P}(\hat{\theta}_n > t)$, et en déduire la densité de $\hat{\theta}_n$.
    \item En déduire la loi de $n(\hat{\theta}_n - \theta)$.
    \item Calculer le risque quadratique de $\hat{\theta}_n$.
\end{enumerate}
\end{exercice}
\end{mdframed}

\smallskip

\begin{mdframed}
\begin{exercice}
Soit $X$ une v.a. symétrique de variance $\sigma^2$ et de fonction de répartition $F$. Soit $(X_i)_{i \ge 1}$ une suite de v.a. i.i.d. de fonction de répartition donnée pour tout $x \in \mathbb{R}$ par $F_\theta(x) = F(x - \theta)$. C'est-à-dire $X_1$ a même loi que $X + \theta$. On s'intéresse à l'estimation de $\theta$.
\begin{enumerate}
    \item Calculer $\mathbb{E}[X]$ et en déduire $\mathbb{E}_\theta[X_1]$.
    \item Dans chacun des cas suivants, calculer l'EMV de $\theta$ :
    \begin{enumerate}
        \item[i--] $F$ est la fonction de répartition de $\mathcal{N}(0, \sigma^2)$.
        \item[ii--] La densité de $X$ est donnée par $f(x) = \frac{3}{4}(1 - x^2)\mathbf{1}_{[-1, 1]}(x)$.
        \item[iii--] La densité de $X$ est donnée sur $\mathbb{R}$ par $f(x) = \frac{1}{2} e^{-|x|}$.
    \end{enumerate}
\end{enumerate}
\end{exercice}
\end{mdframed}

\smallskip

\begin{mdframed}
\begin{exercice}
Soit $X$ une v.a. de loi $\mathbb{P}_\theta$ avec $\theta \in \, ]0, 1[$. Pour tout $\theta \in \, ]0, 1[$, $\mathbb{P}_\theta$ est définie pour tout $k \in \mathbb{N}$, par
\begin{equation*}
\mathbb{P}_\theta(k) = (k + 1)(1 - \theta)^2 \theta^k.
\end{equation*}
On donne
\begin{equation*}
\mathbb{E}_\theta[X] = \frac{2\theta}{1-\theta} \quad \text{et} \quad \text{Var}_\theta(X) = \frac{2\theta}{(1 - \theta)^2}.
\end{equation*}
Soit $X_1, \dots, X_n$ un $n$-échantillon de même loi que $X$.
\begin{enumerate}
    \item Donner un estimateur $\hat{\theta}_n$ de $\theta$ par la méthode des moments.
    \item L'estimateur du maximum de vraisemblance $\tilde{\theta}_n$ de $\theta$ est-il bien défini ?
    \item Étudier la consistance de $\hat{\theta}_n$ et déterminer sa loi limite.
\end{enumerate}
\end{exercice}
\end{mdframed}

\smallskip

\begin{mdframed}
\begin{exercice}
\begin{enumerate}
    \item On définit pour tout $\theta \in \mathbb{R}$, la fonction $f_\theta$ par
    \begin{equation*}
    f_\theta(x) = \begin{cases}
    1 - \theta & \text{si } x \in \, ]{-1/2}, 0], \\
    1 + \theta & \text{si } x \in \, ]0, 1/2], \\
    0 & \text{sinon},
    \end{cases}
    \end{equation*}
    où $\theta$ est un paramètre réel. À quelle condition sur $\theta$, $f_\theta$ est-elle bien une densité sur $\mathbb{R}$ ? On suppose cette condition vérifiée dans la suite.\\
    Soit $X$ une v.a. de loi à densité donnée par $f_\theta$. Soit $(X_1, \dots, X_n)$ un $n$-échantillon de même loi que $X$.
    
    \item Déterminer l'estimateur du maximum de vraisemblance $\hat{\theta}_n$.\\
    \textit{Indication : On pourra faire intervenir la variable aléatoire $Y_n$ égale au nombre de réalisations comprises entre 0 et $1/2$ dans l'échantillon $(X_1, \dots, X_n)$.}
    
    \item $\hat{\theta}_n$ est-il sans biais ? est-il consistant ?
    \item Quelle est la loi limite de $\sqrt{n}(\hat{\theta}_n - \theta)$ ?
\end{enumerate}
\end{exercice}
\end{mdframed}

\smallskip

\begin{mdframed}
\begin{exercice}
Soit $X$ une v.a. de loi $\mathbb{P}_\theta$ avec $\theta \in \, ]0, 1[$. Pour tout $\theta \in \, ]0, 1[$, $\mathbb{P}_\theta$ est définie par
\begin{equation*}
\mathbb{P}_\theta(0) = 1 - \theta, \quad \text{et} \quad \mathbb{P}_\theta(1) = \mathbb{P}_\theta(-1) = \frac{\theta}{2}.
\end{equation*}
\begin{enumerate}
    \item Trouver une expression simple de $\mathbb{P}_\theta(k)$.\\
    \textit{Indication : Penser aux lois de Bernoulli.}
    \item Soit $(X_1, \dots, X_n)$ un $n$-échantillon de même loi que $X$. Déterminer l'EMV $\hat{\theta}_n$ de $\theta$. On note respectivement $T$, $U$ et $V$ le nombre de $X_i$ égales à 0, 1 et $-1$.
    \item $\hat{\theta}_n$ est-il sans biais ?
\end{enumerate}
\end{exercice}
\end{mdframed}

\smallskip

\begin{mdframed}
\begin{exercice}
La hauteur maximale $H$ de la crue annuelle d'un fleuve est observée. Une crue supérieure à 6 mètres serait catastrophique. On modélise $H$ par une variable de Rayleigh : $H$ a une densité donnée par
\begin{equation*}
f(x) = \frac{x}{a} e^{-\frac{x^2}{2a}} \mathbf{1}_{\mathbb{R}_+}(x)
\end{equation*}
où $a > 0$ est un paramètre inconnu. Durant une période de 8 ans, on a observé les hauteurs de crue suivantes en mètres :
\begin{center}
\begin{tabular}{cccccccc}
2.5 & 1.8 & 2.9 & 0.9 & 2.1 & 1.7 & 2.2 & 2.8
\end{tabular}
\end{center}

\begin{enumerate}
    \item Donner l'estimateur du maximum de vraisemblance, $\hat{a}_n$, de $a$.
    \item Montrer que si $X$ suit une loi de Rayleigh de paramètre $a$, alors $X^2$ suit une loi exponentielle de paramètre $1/(2a)$.
    \item Cet estimateur est-il sans biais ? Asymptotiquement normal ?
    \item Une compagnie d'assurance estime qu'une catastrophe n'arrive qu'au plus une fois tous les mille ans. Ceci peut-il être justifié par les observations ?
\end{enumerate}
\end{exercice}
\end{mdframed}
\chapter*{TD 4 - Méthodes d'estimation - 3ème partie}

\setcounter{exercice}{0}

\begin{mdframed}
\begin{exercice}
Reprendre l’Exercice 4 du TD2. Soit $X$ une v.a. symétrique de variance $\sigma^2$ et de fonction de répartition $F$. Soit $(X_i)_{i \ge 1}$ une suite de v.a. i.i.d. de fonction de répartition donnée pour tout $x \in \mathbb{R}$ par $F_{\theta_0}(x) = F(x - \theta_0)$. C’est-à-dire $X_1$ a même loi que $X + \theta_0$. On s’intéresse à l’estimation de $\theta_0$.

\begin{enumerate}
    \item Montrer que si $F$ est continue et strictement croissante sur un voisinage de 0, alors sa médiane vaut 0. En déduire que $\theta = F_\theta^{-1}(1/2)$.
    \item Reprendre la question 2 de l’Exercice 4 du TD2 et calculer les estimateurs de $\theta$ à partir de la médiane empirique dans chacun des 3 cas. Comparez les variances des lois limites entre ces estimateurs et ceux des EMV.
\end{enumerate}
\end{exercice}
\end{mdframed}

\smallskip

\begin{mdframed}
\begin{exercice}
Un statisticien observe $n$ variables aléatoires indépendantes et identiquement distribuées, de loi $\mathcal{U}_{[0,\theta]}$, et veut estimer $\theta > 0$. Il propose les trois estimateurs suivants :
\begin{equation*}
\hat{\theta}_1 = 2\bar{X}_n, \quad \hat{\theta}_2 = 2\hat{F}_n^{-1}(1/2), \quad \hat{\theta}_3 = \sup_{1 \le i \le n} X_i,
\end{equation*}
où $\hat{F}_n^{-1}$ est l’inverse (généralisée) de la fonction de répartition empirique.

\begin{enumerate}
    \item Expliquer les idées conduisant à proposer chacun de ces estimateurs.
    \item Donner la loi limite de $a_{i,n} (\hat{\theta}_i - \theta)$, $i = 1, 2, 3$, où les suites $a_{i,n}$ sont choisies de sorte que l’on obtienne des lois limites non dégénérées.
    \item Quel(s) estimateur(s) préférez-vous ?
\end{enumerate}
\end{exercice}
\end{mdframed}

\smallskip

\begin{mdframed}
\begin{exercice}
On observe $n$ variables aléatoires réelles i.i.d. $X_i$, suivant la loi de Cauchy de centrage $m$ dont la densité est donnée par :
\begin{equation*}
f(x) = \frac{1}{\pi} \frac{1}{1 + (x - m)^2}.
\end{equation*}

\begin{enumerate}
    \item Un statisticien propose la moyenne empirique $\bar{X}_n$ comme estimateur de $m$.
    \begin{enumerate}
        \item[i--] Quelle est la loi de $\bar{X}_n$ ?
        \item[ii--] Étudier la convergence de $\bar{X}_n$ en moyenne quadratique, en probabilité et en loi.
        \item[iii--] Que pensez-vous de cet estimateur ?
    \end{enumerate}
    
    \item Un autre statisticien a l’idée suivante. Il classe les observations par valeurs croissantes et note $X_{(1)} < X_{(2)} < \dots < X_{(n)}$ les valeurs ainsi obtenues. Pour estimer $m$, il propose :
    \begin{equation*}
    \hat{m}_n = \begin{cases}
    X_{(k+1)} & \text{si } n = 2k + 1, \\
    \frac{1}{2}(X_{(k)} + X_{(k+1)}) & \text{si } n = 2k.
    \end{cases}
    \end{equation*}
    \begin{enumerate}
        \item[i--] Montrer que $F_n(\hat{m}_n) \xrightarrow[n \to +\infty]{\text{p.s.}} 1/2$, où $F_n$ est la fonction de répartition empirique.
        \item[ii--] Montrer que, pour la loi considérée, $\mathbb{P}(X_1 < m) = 1/2$. Expliquez l’idée du statisticien.
        \item[iii--] À l’aide du théorème de Glivenko-Cantelli, montrer que $\hat{m}_n \xrightarrow[n \to +\infty]{\text{p.s.}} m$. Commenter.\\
        \textit{Indication : on pourra démontrer que $F(\hat{m}_n) - F(m) \to 0$ p.s. où $F$ désigne la fonction de répartition de $X_1$.}
    \end{enumerate}
\end{enumerate}
\end{exercice}
\end{mdframed}

\smallskip

\begin{mdframed}
\begin{exercice}
Soit $(X_i)_{i \ge 1}$ une suite de v.a. i.i.d. de loi gamma $\Gamma(p, \theta)$, où $p > 0$ et $\theta > 0$. On rappelle si $X \sim \Gamma(p, \theta)$ alors la densité de $X$ est donnée par
\begin{equation*}
f(x) = \frac{\theta^p}{\Gamma(p)} x^{p-1} \exp(-\theta x) \mathbf{1}_{[0,+\infty[}(x),
\end{equation*}
où pour tout $\alpha > 0$,
\begin{equation*}
\Gamma(a) = \int_0^{+\infty} x^{\alpha-1} \exp(-x) dx.
\end{equation*}
On rappelle également que $\mathbb{E}[X] = p/\theta$ et $\text{Var}(X) = p/\theta^2$.

\begin{enumerate}
    \item On suppose $p$ connu et $\theta$ inconnu. Proposer un estimateur consistant de $\theta$.
    \item On suppose $\theta$ connu et $p$ inconnu. Proposer un estimateur consistant de $p$.
    \item On suppose $\theta$ et $p$ inconnus. Proposer des estimateurs consistants de $\theta$ et $p$.
\end{enumerate}
\end{exercice}
\end{mdframed}

\smallskip

\begin{mdframed}
\begin{exercice}
Soit $\theta > 0$ un paramètre inconnu. On considère la densité $f_\theta$ définie sur $\mathbb{R}$ par
\begin{equation*}
f_\theta(x) = \frac{1}{2\theta} \left( \mathbf{1}_{[0,\theta]}(x) + \mathbf{1}_{[2\theta,3\theta]}(x) \right).
\end{equation*}
Soit $(X_i)_{i \ge 1}$ une suite de v.a. i.i.d. de loi de densité $f_\theta$.

\begin{enumerate}
    \item Calculer $\mathbb{E}_\theta[X_1]$ et $\text{Var}_\theta(X_1)$. En déduire un estimateur $\hat{\theta}_n^{MM}$ de $\theta$ par la méthode des moments. Préciser la loi limite de $\hat{\theta}_n^{MM} - \theta$ après une renormalisation convenable.
    \item Calculer la fonction de répartition $F_\theta$ de $X_1$. Représenter $f_\theta$ et $F_\theta$.
    \item Que vaut $q = x_{1/4} = F_\theta^{-1}(1/4)$ ? En déduire un estimateur $\hat{\theta}_n^q$ consistant pour $\theta$ et préciser la loi limite de $\hat{\theta}_n^q - \theta$ après une renormalisation convenable.
    \item Que vaut $Q = x_{3/4} = F_\theta^{-1}(3/4)$ ? En déduire un estimateur $\hat{\theta}_n^Q$ consistant pour $\theta$ et préciser la loi limite de $\hat{\theta}_n^Q - \theta$ après une renormalisation convenable.
    \item Quels estimateurs choisiriez-vous asymptotiquement ?
    \item Soit $X_{(n)} = \max_{1 \le i \le n} X_i$. Pour $0 < x < \theta$, que vaut $\mathbb{P}_\theta(X_{(n)} \le 3\theta - x)$ ? En déduire un estimateur $\hat{\theta}_n^M$ consistant pour $\theta$.
    \item Montrer que $\mathbb{P}_\theta(\hat{\theta}_n^M \le \theta - t/n)$ converge pour tout $t \ge 0$ vers une limite non nulle que l'on calculera. Préciser la loi limite de $n(\theta - \hat{\theta}_n^M)$.
    \item Toujours asymptotiquement, quels estimateurs choisiriez-vous ?
    \item Que vaut $m = x_{1/2} = F_\theta^{-1}(1/2)$ ? Que dire de la médiane empirique comme estimateur de $\theta$ ?
\end{enumerate}
\end{exercice}
\end{mdframed}

\smallskip
\chapter*{TD 5 - Exercices Complet}

\setcounter{exercice}{0}

\begin{mdframed}
\begin{exercice}
On considère une suite de variables aléatoires $(X_i)_{i \ge 1}$ indépendantes et identiquement distribuées de densité $f_\theta$ définie sur $\mathbb{R}$ par
\begin{equation*}
f_{\theta_0}(x) = C(\theta_0)(x - \theta_0)^2 \mathbf{1}_{[\theta_0, \theta_0 + 1]} (x),
\end{equation*}
où $C(\theta_0)$ est une constante (dépendant uniquement de $\theta_0$) et $\theta_0 > 0$ est un paramètre inconnu.

\begin{enumerate}
    \item Montrer que $C(\theta_0) = 3$.
    \item Calculer la fonction de répartition de $X_1$.
    \item Calculer $\mathbb{E}_\theta[X_1 - \theta]$ et en déduire que $\mathbb{E}_\theta[X_1] = \theta + 3/4$.
    \item Calculer $\mathbb{E}_\theta[(X_1 - \theta)^2]$ et en déduire que $\sigma^2 = \text{Var}_\theta(X_1)$ ne dépend pas de $\theta$.
    \item 
    \begin{enumerate}
        \item[i--] À partir de la question 3, définir un estimateur $\hat{\theta}_n$ de $\theta$ par la méthode des moments.
        \item[ii--] $\hat{\theta}_n$ est-il sans biais ? consistant ?
        \item[iii--] Calculer $R(\hat{\theta}_n, \theta)$, le risque quadratique de $\hat{\theta}_n$.
        \item[iv--] Justifier que $\sqrt{n}(\hat{\theta}_n - \theta)$ converge en loi vers une limite que l'on précisera.
        \item[v--] Soit $\alpha \in \, ]0, 1[$ fixé et $t_\alpha = q_{1-\alpha/2}$ le quantile d'ordre $1 - \alpha/2$ de la loi $\mathcal{N}(0, 1)$. Donner la limite de
        \begin{equation*}
        \mathbb{P}_\theta \left( \theta \in \left[ \hat{\theta}_n - \frac{\sigma}{\sqrt{n}}t_\alpha \, ; \, \hat{\theta}_n + \frac{\sigma}{\sqrt{n}}t_\alpha \right] \right).
        \end{equation*}
    \end{enumerate}
    
    \item Dans cette question, on étudie l'estimateur $\tilde{\theta}_n = \max(X_1, \dots, X_n) - 1$.
    \begin{enumerate}
        \item[i--] Calculer la fonction de répartition de $\tilde{\theta}_n$. En déduire que la densité de $\tilde{\theta}_n$ est donnée pour tout $x \in \mathbb{R}$, par
        \begin{equation*}
        f_{\tilde{\theta}_n}(x) = 3n(x + 1 - \theta)^{3n-1} \mathbf{1}_{[\theta-1, \theta]} (x).
        \end{equation*}
        \item[ii--] Montrer que $\mathbb{E}_\theta[\tilde{\theta}_n] = \theta - \frac{1}{3n+1}$. L'estimateur est-il sans biais ?
        \item[iii--] Calculer le risque quadratique de $\tilde{\theta}_n$. En déduire que $\tilde{\theta}_n$ est consistant.
        \item[iv--] Justifier que $\varepsilon_n := 3n(\theta - \tilde{\theta}_n)$ converge en loi vers une limite que l'on précisera. On étudiera pour cela la convergence de la fonction de répartition de $\varepsilon_n$ quand $n \to \infty$.
    \end{enumerate}
    \item Des deux estimateurs $\hat{\theta}_n$ et $\tilde{\theta}_n$ lequel est préférable ? Justifier.
\end{enumerate}
\end{exercice}
\end{mdframed}

\smallskip

\begin{mdframed}
\begin{exercice}
On considère une suite $(U_i)_{i \ge 1}$ de variables aléatoires indépendantes et identiquement distribuées de loi uniforme sur $[0, 1]$. Pour tout $n \ge 1$, on pose
\begin{equation*}
V_n = (U_1 \cdots U_n)^{1/n}.
\end{equation*}

\begin{enumerate}
    \item Montrer qu'avec probabilité 1, $V_n > 0$.
    \item Montrer que $(V_n)_{n \ge 1}$ converge presque sûrement vers une constante qu'on déterminera.\\
    \textit{Indication : on pourra étudier dans un premier temps la convergence de la suite $W_n = \ln(V_n)$, $n \ge 1$.}
    \item On considère maintenant la suite de variables aléatoires $X_n = (e V_n)^{\sqrt{n}}$, $n \ge 1$. Montrer que $X_n$ converge en loi vers une variable aléatoire qu'on déterminera.\\
    \textit{Indication : on considérera $\ln(X_n)$ et on utilisera $\int_0^1 (\ln u)^2 du = 2$.}
\end{enumerate}
\end{exercice}
\end{mdframed}

\smallskip

\begin{mdframed}
\begin{exercice}
On observe $X_1, \dots, X_n$, un $n$-échantillon (v.a. indépendantes et identiquement distribuées) de loi de Poisson de paramètre $\lambda > 0$ inconnu. On rappelle que pour tout $k \in \mathbb{N}$,
\begin{equation*}
\mathbb{P}_\lambda(X_1 = k) = \frac{\lambda^k}{k!} e^{-\lambda}.
\end{equation*}
Le but de cet exercice est d'étudier deux estimateurs du paramètre $p = e^{-\lambda}$.

\begin{enumerate}
    \item On définit pour tout $n \ge 1$, $Y_n = \mathbf{1}_{\{X_n = 0\}}$, et on pose $\bar{Y}_n = \frac{1}{n}\sum_{k=1}^n Y_k$.
    \begin{enumerate}
        \item[i--] Justifier que $\bar{Y}_n$ est un estimateur sans biais et convergeant de $p$.
        \item[ii--] Montrer que $\text{Var}_\lambda(\bar{Y}_n) = \frac{p(1 - p)}{n}$.
    \end{enumerate}
    
    \item On pose $S_n = X_1 + \dots + X_n$, et on considère $N_n = \left( 1 - \frac{1}{n} \right)^{S_n}$.
    \begin{enumerate}
        \item[i--] Montrer que $N_n$ est un estimateur sans biais de $p$.
        \item[ii--] Montrer que $\text{Var}_\lambda(N_n) = e^{-2\lambda}\left( e^{\lambda/n} - 1 \right)$. En déduire que l'estimateur $N_n$ est consistant.
    \end{enumerate}
    
    \item Comparer les risques quadratiques de $\bar{Y}_n$ et $N_n$. Quel estimateur choisiriez-vous ?\\
    On pourra utiliser pour cela l'inégalité suivante : pour tout $\lambda \in \mathbb{R}$, pour tout $n \in \mathbb{N}$,
    \begin{equation}
    n(e^{\lambda/n} - 1) \le e^\lambda - 1.
    \end{equation}
    \textit{Question facultative : démontrer l'inégalité.}
    
    \item On étudie dans cette question le caractère asymptotiquement gaussien de $N_n$.
    \begin{enumerate}
        \item[i--] On rappelle que $\mathbb{E}_\lambda[X_1] = \text{Var}_\lambda(X_1) = \lambda$. Justifier que
        \begin{equation*}
        \sqrt{n}\left( \frac{S_n}{n} - \lambda \right) \xrightarrow[n \to +\infty]{\text{loi}} \sqrt{\lambda}Z
        \end{equation*}
        avec $Z \sim \mathcal{N}(0, 1)$.
        \item[ii--] On pose $\alpha_n = n \ln(1 - 1/n)$. On rappelle le développement asymptotique suivant lorsque $n \to +\infty$ :
        \begin{equation*}
        \alpha_n = -1 - \frac{1}{2n} + o\left(\frac{1}{n}\right).
        \end{equation*}
        En déduire que
        \begin{equation*}
        \sqrt{n}\left( \frac{S_n}{n}\alpha_n + \lambda \right) \xrightarrow[n \to +\infty]{\text{loi}} -\sqrt{\lambda}Z.
        \end{equation*}
        \textit{Indication : on fera intervenir la quantité $\sqrt{n}\left( \frac{S_n}{n} - \lambda \right)\alpha_n$.}
        \item[iii--] Montrer que
        \begin{equation*}
        \sqrt{n}(N_n - p) \xrightarrow[n \to +\infty]{\text{loi}} e^{-\lambda}\sqrt{\lambda}Z.
        \end{equation*}
        \item[iv--] Soit $\alpha > 0$. En déduire un intervalle $I$ de la forme $[N_n \pm U_n]$ tel que $U_n$ dépend uniquement de $n$, $N_n$ et $\alpha$ vérifiant
        \begin{equation*}
        \mathbb{P}_\lambda(p \in I) \xrightarrow[n \to +\infty]{} 1 - \alpha.
        \end{equation*}
    \end{enumerate}
\end{enumerate}
\end{exercice}
\end{mdframed}

\smallskip

\begin{mdframed}
\begin{exercice}
On observe $X_1, \dots, X_n$, un $n$-échantillon de densité donnée pour tout $x \in \mathbb{R}$, par
\begin{equation*}
f_\theta(x) = \frac{1}{\theta}(1 - x)^{\frac{1}{\theta} - 1} \mathbf{1}_{]0, 1[}(x),
\end{equation*}
où $\theta > 0$ est un paramètre inconnu.

\begin{enumerate}
    \item Vérifier que, pour tout $\theta > 0$, $f_\theta$ est bien une densité sur $]0, 1[$.
    \item Calculer la fonction de répartition de $X_1$.
    \item 
    \begin{enumerate}
        \item[i--] Calculer $\mathbb{E}_\theta[X_1]$. En déduire un estimateur $\hat{\theta}_n$ de $\theta$ par la méthode des moments.\\
        On pourra utiliser l'égalité $\mathbb{E}_\theta[X_1] = \int_0^{+\infty} \mathbb{P}_\theta(X_1 > t)dt$.
        \item[ii--] Justifier que $\hat{\theta}_n$ est consistant.
        \item[iii--] Justifier que $\hat{\theta}_n$ est biaisé (on ne demande pas de calculer son biais).\\
        \textit{Indication : on pourra utiliser l'inégalité de Jensen.}
    \end{enumerate}
    \item On définit, pour tout $i \ge 1$, la variable $Y_i = -\log(1 - X_i)$. Montrer que $Y_1$ suit la loi exponentielle de paramètre $1/\theta$.
    \item 
    \begin{enumerate}
        \item[i--] Montrer que $\tilde{\theta}_n = \bar{Y}_n$ est l'estimateur du maximum de vraisemblance du paramètre $\theta$.
        \item[ii--] Justifier que $\tilde{\theta}_n$ est consistant.
        \item[iii--] Montrer que $\tilde{\theta}_n$ est sans biais et calculer son risque quadratique.
        \item[iv--] Déterminer la limite en loi de $\sqrt{n}(\tilde{\theta}_n - \theta)$ quand $n \to +\infty$.
    \end{enumerate}
\end{enumerate}
\end{exercice}
\end{mdframed}

\chapter{Borne de Cramér-Rao}

L'inégalité de Cramér-Rao fournit une minoration du risque quadratique de tout estimateur sans biais, dans un modèle paramétrique dominé satisfaisant quelques hypothèses de régularité. Elle exprime ainsi une limite intrinsèque au problème d'estimation : aucun estimateur sans biais ne peut faire mieux que la borne, quelle que soit la méthode employée pour le construire.

\textit{Ce résultat ne concerne que les estimateurs sans biais.}

\medskip
Dans toute cette section, on suppose les hypothèses de régularité suivantes, qui définissent ce que l'on appellera un \textbf{modèle régulier}:

\begin{itemize}
  \item[$(\mathcal{H}_0)$] Il existe une mesure positive $\nu$ telle que, pour tout $\theta \in \Theta$, $P_\theta$ admet une densité $f_\theta$ par rapport à $\nu$.
  \item[$(\mathcal{H}_1)$] $\Theta \subset \mathbb{R}$ est un ouvert et, pour tout $x$ de l'espace des observations $\mathcal{X}$, l'application $\theta \longmapsto f_\theta(x)$ est dérivable sur $\Theta$.
  \item[$(\mathcal{H}_2)$] Les supports des densités $f_\theta$, $\theta \in \Theta$, ne dépendent pas de $\theta$ ; on note leur valeur commune $S = \{x \in \mathcal{X} : f_\theta(x) > 0\}$.
  \item[$(\mathcal{H}_3)$] Pour tout $\theta \in \Theta$, la variable aléatoire $\frac{\partial}{\partial\theta} \ln f_\theta(X)$ est centrée sous $P_\theta$.
  \item[$(\mathcal{H}_4)$] Pour tout $\theta \in \Theta$, la variable aléatoire $\frac{\partial}{\partial\theta} \ln f_\theta(X)$ appartient à $L^2(P_\theta)$.
\end{itemize}

\begin{remarque}
L'hypothèse $(\mathcal{H}_3)$ est en réalité une conséquence d'une interversion entre
dérivation et intégration. En effet, comme $\int_S f_\theta \, d\nu = 1$ pour tout
$\theta$, si l'on peut dériver sous le signe intégral on obtient

$$
  \int_S \frac{\partial}{\partial\theta} f_\theta(x)\, d\nu(x) = 0,
  \qquad\text{soit}\qquad
  E_\theta\!\left[ \frac{\partial}{\partial\theta} \ln f_\theta(X) \right]
  = \int_S \frac{\partial_\theta f_\theta}{f_\theta}\, f_\theta \, d\nu = 0.
$$

Dans les modèles usuels, cette interversion est légitime dès que le théorème de dérivation des intégrales à paramètre s'applique (existence d'une domination intégrable de $\partial_\theta f_\theta$ au voisinage de $\theta$).
\end{remarque}

\begin{remarque}
L'hypothèse (H2) est essentielle : c'est elle qui exclut les modèles où le support
dépend du paramètre. La famille exponentielle
$\{f_\theta(x) = \theta\, e^{-\theta x}\,\mathbb{1}_{x \geq 0},\ \theta > 0\}$ a un
support $S = [0, +\infty[$ indépendant de $\theta$ et constitue un modèle régulier.
En revanche, la famille des lois uniformes $\mathcal{U}([0,\theta])$, de densité
$f_\theta(x) = \frac{1}{\theta}\mathbb{1}_{[0,\theta]}(x)$, a un support
$[0,\theta]$ qui dépend de $\theta$ : (H2) n'est pas vérifiée et la borne de
Cramér-Rao ne s'applique pas. C'est d'ailleurs pour ce type de modèle que l'on peut
obtenir des estimateurs convergeant à vitesse $1/n$, plus rapide que la vitesse
usuelle $1/\sqrt{n}$.
\end{remarque}

% ---------------------------------------------------------------------
\begin{mdframed}
\begin{definition}[Information de Fisher]
Dans un modèle paramétrique régulier, l'\textbf{information de Fisher} est la fonction $I : \Theta \to [0, +\infty[$ définie par

$$
  I(\theta)
  = E_\theta\!\left[ \left( \frac{\partial}{\partial\theta} \ln L(\theta) \right)^{2} \right]
  = \operatorname{Var}_\theta\!\left( \frac{\partial}{\partial\theta} \ln L(\theta) \right),
$$

où $L(\theta)$ désigne la vraisemblance. La variable aléatoire
$\dfrac{\partial}{\partial\theta} \ln L(\theta)$ est appelée \textbf{fonction de
score}.
\end{definition}
\end{mdframed}

L'égalité entre les deux expressions de $I(\theta)$ provient directement de (H3) :
le score étant centré, son espérance du carré coïncide avec sa variance.

\begin{remarque}[Cas multidimensionnel]
La définition est donnée pour $\theta \in \mathbb{R}$. Lorsque $\theta$ est un
vecteur de dimension $p > 1$, l'information de Fisher est la matrice $p \times p$
\[
  I(\theta)
  = \left( E_\theta\!\left[ \frac{\partial}{\partial\theta_i} \ln L(\theta)\,
    \frac{\partial}{\partial\theta_j} \ln L(\theta) \right] \right)_{1 \leq i,j \leq p},
\]
c'est-à-dire la matrice de covariance du vecteur score.
\end{remarque}

\begin{proposition}[Expression par la dérivée seconde]
On suppose de plus que, pour tout $x \in S$, l'application
$\theta \mapsto f_\theta(x)$ est de classe $C^2$ sur $\Theta$, et que l'on peut
dériver deux fois sous le signe intégral l'identité $\int_S f_\theta\, d\nu = 1$
(domination intégrable des dérivées première et seconde au voisinage de $\theta$).
Alors
\[
  I(\theta) = -\, E_\theta\!\left[ \frac{\partial^{2}}{\partial\theta^{2}} \ln L(\theta) \right].
\]
\end{proposition}

\begin{proof}
En dérivant deux fois la relation $\int_S f_\theta\, d\nu = 1$, on obtient
$\int_S \partial_\theta^2 f_\theta \, d\nu = 0$. Par ailleurs, partout sur $S$,
\[
  \frac{\partial^{2}}{\partial\theta^{2}} \ln f_\theta
  = \frac{\partial_\theta^{2} f_\theta}{f_\theta}
    - \left( \frac{\partial_\theta f_\theta}{f_\theta} \right)^{2}
  = \frac{\partial_\theta^{2} f_\theta}{f_\theta}
    - \left( \frac{\partial}{\partial\theta} \ln f_\theta \right)^{2}.
\]
En intégrant contre $f_\theta$ sous $P_\theta$ :
\[
  E_\theta\!\left[ \frac{\partial^{2}}{\partial\theta^{2}} \ln f_\theta(X) \right]
  = \int_S \partial_\theta^{2} f_\theta \, d\nu
    - E_\theta\!\left[ \left( \frac{\partial}{\partial\theta} \ln f_\theta(X) \right)^{2} \right]
  = 0 - I(\theta) = -\,I(\theta). \qedhere
\]
\end{proof}

\begin{remarque}
On ne peut pas prendre cette dernière formule comme \emph{définition} de
l'information de Fisher : elle n'est valable que sous les hypothèses de classe
$C^2$ et d'interversion ci-dessus. En revanche, la définition par la variance du
score ne requiert que (H1)--(H4). En pratique, lorsque ces hypothèses sont
satisfaites, la formule par la dérivée seconde est souvent la plus commode pour le
calcul.
\end{remarque}

\begin{remarque}[Interprétation]
L'information de Fisher mesure la quantité d'information que l'observation apporte
sur le paramètre $\theta$. La formule $I(\theta) = -E_\theta[\partial_\theta^{2}
\ln L(\theta)]$ l'identifie à la \emph{courbure moyenne} de la log-vraisemblance au
point $\theta$ : plus la log-vraisemblance est «~pointue~» autour de la vraie
valeur, plus $I(\theta)$ est grande, et plus $\theta$ est facile à estimer (la
borne de Cramér-Rao $\frac{1}{n I(\theta)}$ est alors petite). À l'inverse, une
log-vraisemblance plate, $I(\theta)$ proche de $0$, traduit que des valeurs
voisines de $\theta$ produisent des lois quasi indiscernables : l'estimation est
difficile et la borne explose. C'est l'analogue quantitatif de l'identifiabilité
vue au chapitre~1.
\end{remarque}

\begin{remarque}[Effet d'un reparamétrage]
L'information de Fisher dépend du paramétrage choisi. Si l'on pose $\eta =
\varphi(\theta)$ avec $\varphi$ un $C^1$-difféomorphisme, alors, en dérivant le
score par composition, l'information dans le paramétrage $\eta$ s'écrit
\[
  I_\eta(\eta) = \frac{I_\theta(\theta)}{\bigl(\varphi'(\theta)\bigr)^{2}},
  \qquad \theta = \varphi^{-1}(\eta).
\]
On retrouve la cohérence avec la méthode delta (\S1.4) : la borne $\frac{1}{n
I_\eta(\eta)} = \frac{(\varphi'(\theta))^{2}}{n I_\theta(\theta)}$ se transforme
exactement comme la variance asymptotique sous une transformation régulière.
\end{remarque}

% ---------------------------------------------------------------------
\begin{mdframed}
\begin{proposition}[Information de Fisher d'un $n$-échantillon]
Si le modèle $\{P_\theta, \theta \in \Theta\}$ est régulier, d'information de
Fisher $I(\theta)$, alors le modèle produit
$\{P_\theta^{\otimes n}, \theta \in \Theta\}$ est régulier, d'information de Fisher
\[
  I_n(\theta) = n\, I(\theta).
\]
\end{proposition}
\end{mdframed}

\begin{proof}
Pour un $n$-échantillon $X = (X_1, \dots, X_n)$ i.i.d., la vraisemblance se
factorise : $L_n(\theta) = \prod_{i=1}^{n} f_\theta(X_i)$, de sorte que la
log-vraisemblance, et donc le score, sont des sommes :
\[
  \frac{\partial}{\partial\theta} \ln L_n(\theta)
  = \sum_{i=1}^{n} \frac{\partial}{\partial\theta} \ln f_\theta(X_i).
\]
Les variables $\frac{\partial}{\partial\theta} \ln f_\theta(X_i)$ sont
indépendantes, identiquement distribuées, centrées (par (H3)) et de variance
$I(\theta)$ chacune. La variance d'une somme de variables indépendantes étant la
somme des variances,
\[
  I_n(\theta)
  = \operatorname{Var}_\theta\!\left( \sum_{i=1}^{n}
    \frac{\partial}{\partial\theta} \ln f_\theta(X_i) \right)
  = n\, \operatorname{Var}_\theta\!\left( \frac{\partial}{\partial\theta} \ln f_\theta(X_1) \right)
  = n\, I(\theta). \qedhere
\]
\end{proof}

% ---------------------------------------------------------------------
On considère désormais une quantité d'intérêt $g(\theta)$, où $g : \Theta \to
\mathbb{R}$ est dérivable, et un estimateur $T = T(X)$ \emph{sans biais} de
$g(\theta)$. On ajoute l'hypothèse de régularité suivante.

\begin{itemize}
  \item[\textbf{(H5)}] Pour tout $\theta \in \Theta$, on peut dériver sous le signe
        intégral :
        \[
          \frac{\partial}{\partial\theta} \int_S T(x)\, f_\theta(x)\, d\nu(x)
          = \int_S T(x)\, \frac{\partial}{\partial\theta} f_\theta(x)\, d\nu(x).
        \]
\end{itemize}

Cette écriture vaut aussi bien pour les modèles discrets que pour les modèles
dominés par la mesure de Lebesgue.

\begin{mdframed}
\begin{proposition}[Inégalité de Cramér-Rao]
Sous les hypothèses (H1)--(H4) d'un modèle régulier, si $I(\theta) > 0$ pour tout
$\theta \in \Theta$, alors tout estimateur $T_n = T(X_1, \dots, X_n)$ sans biais du
paramètre réel $g(\theta)$, vérifiant (H5) et de moment d'ordre $2$ fini
($E_\theta[T_n^2] < +\infty$), satisfait
\[
  \forall\, \theta \in \Theta, \qquad
  E_\theta\!\left[ (T_n - g(\theta))^{2} \right]
  = \operatorname{Var}_\theta(T_n)
  \geq \frac{\bigl(g'(\theta)\bigr)^{2}}{I_n(\theta)}
  = \frac{\bigl(g'(\theta)\bigr)^{2}}{n\, I(\theta)}.
\]
Le membre de droite est appelé \textbf{borne de Cramér-Rao}.
\end{proposition}
\end{mdframed}

\begin{proof}
On raisonne au niveau de l'observation $X$ (le passage au $n$-échantillon se fait
ensuite en remplaçant $I(\theta)$ par $I_n(\theta) = n I(\theta)$). Comme $T$ est
sans biais, $E_\theta[T(X)] = g(\theta)$, soit
$\int_S T(x) f_\theta(x)\, d\nu(x) = g(\theta)$. En dérivant en $\theta$ et en
utilisant (H5),
\[
  g'(\theta)
  = \int_S T(x)\, \frac{\partial}{\partial\theta} f_\theta(x)\, d\nu(x)
  = \int_S T(x)\, \frac{\partial_\theta f_\theta(x)}{f_\theta(x)}\, f_\theta(x)\, d\nu(x)
  = E_\theta\!\left[ T(X)\, \frac{\partial}{\partial\theta} \ln f_\theta(X) \right].
\]
Le score étant centré (H3), on peut soustraire $g(\theta)$ sans rien changer :
\[
  g'(\theta)
  = E_\theta\!\left[ \bigl(T(X) - g(\theta)\bigr)\, \frac{\partial}{\partial\theta} \ln f_\theta(X) \right].
\]
On reconnaît le produit scalaire $\langle h_1, h_2 \rangle_\theta =
E_\theta[h_1(X) h_2(X)]$ entre les deux variables centrées $T(X) - g(\theta)$ et
$\frac{\partial}{\partial\theta} \ln f_\theta(X)$. L'inégalité de Cauchy-Schwarz
donne
\[
  \bigl(g'(\theta)\bigr)^{2}
  \leq E_\theta\!\left[ \bigl(T(X) - g(\theta)\bigr)^{2} \right]
       \cdot E_\theta\!\left[ \left( \frac{\partial}{\partial\theta} \ln f_\theta(X) \right)^{2} \right]
  = \operatorname{Var}_\theta(T)\cdot I(\theta),
\]
la dernière égalité résultant du caractère sans biais de $T$ et de la définition de
$I(\theta)$. Comme $I(\theta) > 0$, on en déduit
\[
  \operatorname{Var}_\theta(T) \geq \frac{\bigl(g'(\theta)\bigr)^{2}}{I(\theta)}.
\]
Pour un $n$-échantillon, le même raisonnement appliqué à $L_n$ remplace $I(\theta)$
par $I_n(\theta) = n I(\theta)$, ce qui donne la borne annoncée. \qedhere
\end{proof}

\begin{remarque}[Version matricielle]
Lorsque $\theta \in \mathbb{R}^{p}$, l'inégalité devient une inégalité entre
matrices : pour tout estimateur sans biais $T$ de $\theta$,
\[
  \operatorname{Var}_\theta(T) \;\succeq\; I_n(\theta)^{-1},
\]
au sens où $\operatorname{Var}_\theta(T) - I_n(\theta)^{-1}$ est symétrique positive
(ordre de Loewner), $I_n(\theta)$ étant la matrice d'information de Fisher. En
particulier, en regardant les termes diagonaux, la variance de chaque composante
$T_j$ est minorée par le $j$-ème terme diagonal de $I_n(\theta)^{-1}$.
\end{remarque}

% ---------------------------------------------------------------------
\begin{mdframed}
\begin{definition}[Estimateur efficace]
Un estimateur $T$ sans biais de $g(\theta)$ qui réalise l'égalité dans l'inégalité
de Cramér-Rao, c'est-à-dire tel que
\[
  \operatorname{Var}_\theta(T) = \frac{\bigl(g'(\theta)\bigr)^{2}}{n\, I(\theta)}
  \qquad \forall\, \theta \in \Theta,
\]
est dit \textbf{efficace} ; on dit aussi qu'il \emph{atteint la borne de
Cramér-Rao}.
\end{definition}
\end{mdframed}

\begin{remarque}
Un estimateur efficace est, par construction, de risque quadratique minimal parmi
les estimateurs sans biais : c'est en particulier un estimateur UVMB (au sens de la
Remarque sur les estimateurs sans biais du chapitre~1). La réciproque est fausse :
un estimateur peut être UVMB sans atteindre la borne de Cramér-Rao, car cette borne
n'est pas toujours atteignable.
\end{remarque}

\begin{remarque}[Cas d'égalité]
L'égalité dans Cauchy-Schwarz a lieu si et seulement si les deux variables sont
$P_\theta$-presque sûrement proportionnelles. Un estimateur sans biais efficace
n'existe donc que si la fonction de score s'écrit
\[
  \frac{\partial}{\partial\theta} \ln L(\theta)
  = A(\theta)\,\bigl(T(X) - g(\theta)\bigr)
  \qquad P_\theta\text{-p.s.,}
\]
pour une certaine fonction $A(\theta)$. Cette contrainte n'est satisfaite que par
les modèles de type \emph{famille exponentielle}, ce qui explique que l'efficacité
soit une propriété rare.
\end{remarque}

% =====================================================================
\subsection*{Exemples}

\begin{exemple}[Modèle de Bernoulli]
Soit $(X_1, \dots, X_n)$ un $n$-échantillon i.i.d. de loi $\mathcal{B}(\theta)$,
$\theta \in {]0,1[}$, et $g(\theta) = \theta$. La densité (par rapport à la mesure
de comptage) est $f_\theta(x) = \theta^{x}(1-\theta)^{1-x}$, $x \in \{0,1\}$, d'où
\[
  \ln f_\theta(x) = x \ln\theta + (1-x)\ln(1-\theta),
  \qquad
  \frac{\partial}{\partial\theta} \ln f_\theta(x)
  = \frac{x}{\theta} - \frac{1-x}{1-\theta}
  = \frac{x - \theta}{\theta(1-\theta)}.
\]
Comme $\operatorname{Var}_\theta(X_1) = \theta(1-\theta)$,
\[
  I(\theta)
  = \operatorname{Var}_\theta\!\left( \frac{X_1 - \theta}{\theta(1-\theta)} \right)
  = \frac{\theta(1-\theta)}{\theta^{2}(1-\theta)^{2}}
  = \frac{1}{\theta(1-\theta)}.
\]
La borne de Cramér-Rao pour un estimateur sans biais de $\theta$ vaut donc
$\dfrac{1}{n I(\theta)} = \dfrac{\theta(1-\theta)}{n}$. Or la moyenne empirique
$\overline{X}_n$ est sans biais et $\operatorname{Var}_\theta(\overline{X}_n) =
\frac{\theta(1-\theta)}{n}$ : elle \textbf{atteint la borne} et est donc efficace.
\end{exemple}

\begin{exemple}[Modèle de Poisson]
Soit $(X_1, \dots, X_n)$ i.i.d. de loi $\mathcal{P}(\theta)$, $\theta > 0$, et
$g(\theta) = \theta$. On a $P_\theta(X_1 = k) = e^{-\theta}\theta^{k}/k!$, d'où
\[
  \ln f_\theta(k) = -\theta + k\ln\theta - \ln(k!),
  \qquad
  \frac{\partial}{\partial\theta} \ln f_\theta(k) = -1 + \frac{k}{\theta}
  = \frac{k - \theta}{\theta}.
\]
Comme $\operatorname{Var}_\theta(X_1) = \theta$,
\[
  I(\theta) = \frac{\operatorname{Var}_\theta(X_1)}{\theta^{2}} = \frac{1}{\theta},
  \qquad\text{borne de Cramér-Rao} = \frac{\theta}{n}.
\]
La moyenne empirique $\overline{X}_n$ est sans biais et de variance
$\frac{\theta}{n}$ : elle est efficace.
\end{exemple}

\begin{exemple}[Modèle exponentiel]
Soit $(X_1, \dots, X_n)$ i.i.d. de loi $\mathcal{E}(\theta)$, de densité
$f_\theta(x) = \theta e^{-\theta x}\mathbb{1}_{x \geq 0}$, $\theta > 0$. Alors
\[
  \ln f_\theta(x) = \ln\theta - \theta x,
  \qquad
  \frac{\partial}{\partial\theta} \ln f_\theta(x) = \frac{1}{\theta} - x,
  \qquad
  I(\theta) = \operatorname{Var}_\theta(X_1) = \frac{1}{\theta^{2}}.
\]
On ne peut pas estimer $\theta$ lui-même sans biais (cf. Exemple~1.4 du
chapitre~1). On estime donc plutôt $g(\theta) = \tfrac{1}{\theta} = E_\theta[X_1]$,
avec $g'(\theta) = -\tfrac{1}{\theta^{2}}$. La borne de Cramér-Rao est
\[
  \frac{\bigl(g'(\theta)\bigr)^{2}}{n I(\theta)}
  = \frac{1/\theta^{4}}{n/\theta^{2}} = \frac{1}{n\theta^{2}}.
\]
La moyenne empirique $\overline{X}_n$ est sans biais pour $\tfrac{1}{\theta}$ et de
variance $\frac{1}{n}\operatorname{Var}_\theta(X_1) = \frac{1}{n\theta^{2}}$ : elle
est efficace pour $g(\theta) = \tfrac{1}{\theta}$.
\end{exemple}

\begin{exemple}[Modèle gaussien]
Soit $(X_1, \dots, X_n)$ i.i.d. de loi $\mathcal{N}(\theta, \sigma^{2})$, avec
$\sigma^{2} > 0$ \emph{connue} et $g(\theta) = \theta$. La densité est
$f_\theta(x) = \frac{1}{\sqrt{2\pi\sigma^{2}}}\exp\!\bigl(-\frac{(x-\theta)^{2}}{2\sigma^{2}}\bigr)$,
d'où
\[
  \frac{\partial}{\partial\theta} \ln f_\theta(x) = \frac{x - \theta}{\sigma^{2}},
  \qquad
  I(\theta) = \frac{\operatorname{Var}_\theta(X_1)}{\sigma^{4}}
  = \frac{\sigma^{2}}{\sigma^{4}} = \frac{1}{\sigma^{2}}.
\]
La borne de Cramér-Rao vaut $\dfrac{\sigma^{2}}{n}$. La moyenne empirique
$\overline{X}_n$ est sans biais et de variance $\frac{\sigma^{2}}{n}$ : elle est
efficace. On retrouve ainsi que l'estimateur $\overline{X}_n$ de l'Exemple~1.3
(cas $\sigma^{2} = 1$, de risque $1/n$) n'est pas seulement raisonnable : il est
optimal au sens de Cramér-Rao.
\end{exemple} 

\chapter*{TD 6 - Optimalité}

\section*{Rappels}
La condition d’échange entre intégrale et dérivation est typiquement garantie par l’existence pour tout $\theta$ d’un voisinage ouvert $V$ tel que $T sup_{\theta^\prime V} \frac{\partial}{\partial \theta} f_{\theta^\prime}$ soit intégrable ; le théorème de dérivation des intégrales à paramètres assure alorsla que la condition est vérifiée.

Dans les TD, la borne de Cramer Rao et information de Fisher ne seront évoquées que si cette condition
est implicitement satisfaite.

\medskip

\setcounter{exercice}{0}

\begin{mdframed}
\begin{exercice}
Soit $X_1, \dots, X_n$ un $n$-échantillon de densité 

$$
f(x) = \frac{1}{\theta} \exp\left(-\frac{x}{\theta}\right) \mathbb{1}_{[0,+\infty[}(x)
$$

\begin{enumerate}
    \item Montrer qe l'estimateur du maximum de vraisemblance de $\theta$ est $\hat{\theta} = \bar{X}$
    \item Déterminer l'information de Fisher du modèle.
    \item L'estimateur $\hat{\theta}$ est-il efficace pour estimer $\theta$ ?
\end{enumerate}
\end{exercice}
\end{mdframed}

\medskip

\begin{mdframed}
\begin{exercice}
Soit $(X_1, \dots, X_n)$ un $n$-échantillon de densité 

$$
f(x) = 2\theta x e^{-\theta x^2} \mathbb{1}_{x>0},
$$

où $\theta > 0$. On pourra admettre que si $Z \sim \Gamma(\alpha, \beta)$, alors :

\begin{itemize}
\item $\mathbb{E}[Z] = \frac{\alpha}{\beta}$, 
\item $\mathbb{E}[Z^2] = \frac{\alpha(\alpha+1)}{\beta^2}$, 
\item $\mathbb{E}\left[\frac{1}{Z}\right] = \frac{\beta}{\alpha-1}$, 
\item $\mathbb{E}\left[\frac{1}{Z^2}\right] = \frac{\beta^2}{(\alpha-1)(\alpha-2)}$
\end{itemize}

(propriétés à savoir retrouver).
\begin{enumerate}
    \item Calculer l'estimateur du maximum de vraisemblance $\hat{\theta}$ de $\theta$.
    \item Montrer que $Y_i = \theta X_i^2$ suit la loi exponentielle de paramètre $1$. En déduire la loi de $Z = \sum_{i=1}^n \theta X_i^2$.
\item Retrouver les propriétés de $Z$.
\item Déterminer l'information de Fisher du modèle.
    \item L'EMV est-il biaisé ? En déduire un estimateur $\hat{\theta}_1$ non biaisé.
    \item Calculer la variance de $\hat{\theta}_1$. L'estimateur $\hat{\theta}_1$ est-il efficace pour estimer $\theta$ ?
\end{enumerate}
\end{exercice}
\end{mdframed}

\medskip

\begin{mdframed}
\begin{exercice}
Soit $(X_1, \dots, X_n)$ $n$ variables i.i.d. suivant une loi de Poisson de paramètre $\theta$. On rappelle les propriétés suivantes de cette loi :

\begin{itemize}
    \item Sa fonction génératrice des moments est $M(t) = \exp(\theta(e^t - 1))$ pour $t \in \mathbb{R}$.
    \item $\mathbb{E}[X] = \theta$ ; $\mathbb{E}[X^2] = \theta^2 + \theta$ ; $\mathbb{E}[X^3] = \theta^3 + 3\theta^2 + \theta$ ; $\mathbb{E}[X^4] = \theta^4 + 6\theta^3 + 7\theta^2 + \theta$.
    \item $\sum_{i=1}^n X_i$ suit la loi de Poisson de paramètre $n\theta$.
\end{itemize}

On rappelle également :

\begin{itemize}
    \item L'estimateur du maximum de vraisemblance de $\theta$ est la moyenne empirique $\bar{X}$.
    \item L'information de Fisher du modèle vaut $I_n(\theta) = \frac{n}{\theta}$.
    \item L'estimateur $\bar{X}$ atteint la borne de Cramér--Rao : $\text{Var}(\bar{X}) = \frac{\theta}{n} = \frac{1}{I_n(\theta)}$ et est donc efficace pour estimer $\theta$.
\end{itemize}

\begin{enumerate}
    \item Dans cette question, on cherche à estimer $g(\theta) = e^{-\theta}$. On pose $Y_i = \mathbb{1}_{\{X_i=0\}}$ et $\bar{Y} = \frac{1}{n}\sum_{i=1}^n Y_i$.
    \begin{enumerate}
        \item Vérifier que $\bar{Y}$ est un estimateur sans biais de $g(\theta)$.
        \item Calculer la borne de Cramér--Rao pour $\bar{Y}$.
        \item Cet estimateur est-il efficace pour estimer $g(\theta)$ ?
    \end{enumerate}
    \item Dans cette question, on cherche à estimer $h(\theta) = \theta^2$. On pose $T_n = \bar{X}^2 - \frac{1}{n}\bar{X}$.
    \begin{enumerate}
        \item Vérifiez que $T_n$ est un estimateur sans biais de $h(\theta)$.
        \item Calculer la borne de Cramér--Rao pour $T_n$.
        \item Cet estimateur est-il efficace pour estimer $h(\theta)$ ?
    \end{enumerate}
\end{enumerate}
\end{exercice}
\end{mdframed}






























\medskip

\begin{mdframed}
\begin{exercice}
Soit $\theta > 0$ et $(X_1, \dots, X_n)$ un $n$-échantillon de densité :

$$
f_\theta(x) = \begin{cases} 
\frac{1}{\theta} x^{\frac{1}{\theta} - 1} & \text{si } 0 < x < 1, \\ 
0 & \text{sinon.} 
\end{cases} 
$$

\begin{enumerate}
    \item Calculer l'estimateur du maximum de vraisemblance $\hat{\theta}$ de $\theta$.
    \item Cet estimateur est-il biaisé ? \\
    \textit{Indication :} déterminer la loi de probabilité de $Y_i = -\frac{1}{\theta} \log X_i$ puis de $Z = \sum_{i=1}^n Y_i$.
    \item Calculer la variance de $\hat{\theta}$.
    \item Montrer que $\ell'(\theta) = \frac{n}{\theta^2}(\hat{\theta} - \theta)$, où $\ell(\theta)$ désigne la log-vraisemblance des observations.
    \item En déduire que l'estimateur $\hat{\theta}$ atteint la borne de Cramér-Rao.
\end{enumerate}
\end{exercice}
\end{mdframed}

\medskip

\begin{mdframed}
\begin{exercice}
On considère le modèle des lois $P_\theta$ de densité $f_\theta$, pour $\theta > 0$, données par :

$$
f_\theta(x) = \frac{\theta \exp\left(-\theta(\sqrt{x} - 1)\right)}{2\sqrt{x}} \mathbb{1}_{x \ge 1} 
$$

\begin{enumerate}
    \item Si $X$ est de loi $P_\theta$, quelle est la loi de $\sqrt{X}$ ?
    \item Vérifier que la variable aléatoire $\frac{\partial \ln f_\theta(X)}{\partial \theta}$ est centrée sous $P_\theta$.
    \item Calculer l'information de Fisher du modèle.
    \item On considère un $n$-échantillon $X_1, \dots, X_n$ de loi $P_\theta$. Montrer qu'il existe un unique estimateur du maximum de vraisemblance $\hat{\theta}_n$ et le calculer.
    \item Soit $g$ la fonction donnée par $g(u) = \frac{1}{u}$ pour tout $u > 0$. Écrire l'inégalité de Cramér-Rao pour l'estimation de $g(\theta)$.
    \item Montrer qu'il existe un estimateur sans biais efficace de $g(\theta)$.
\end{enumerate}
\end{exercice}
\end{mdframed}

\chapter{Intervalles de confiance}

\section{Définition et interprétation}

On se place dans un modèle statistique $\{P_\theta, \theta \in \Theta\}$ avec $\Theta \subset \mathbb R$. On observe $X$ et on souhaite non plus seulement estimer $\theta$ par une valeur ponctuelle $\hat\theta$, mais \textit{quantifier l'incertitude} autour de cet estimateur.

\medskip

\begin{mdframed}
\begin{definition}[Intervalle de confiance]
Soit $\alpha \in\, ]0,1[$. Un \textit{intervalle de confiance} (IC) de niveau $1-\alpha$ pour $\theta$ est un intervalle aléatoire $I(X) = [\,T_-(X),\, T_+(X)\,]$ tel que

$$
\forall\,\theta \in \Theta, \qquad P_\theta\big(\theta \in I(X)\big) \ge 1-\alpha.
$$

Le réel $1-\alpha$ est le \textit{niveau de confiance}, $\alpha$ le \textit{risque}.
\end{definition}
\end{mdframed}

\begin{remarque}[Interprétation fréquentiste] 
L'interprétation correcte est la suivante : si l'on répète l'expérience un grand nombre de fois et que l'on construit à chaque fois l'intervalle $I(X)$, alors au moins une proportion $1-\alpha$ de ces intervalles contient la vraie valeur $\theta$.

Il est \textit{incorrect} de dire que « $\theta$ appartient à $I(x)$ avec probabilité $1-\alpha$ » une fois l'observation $x$ fixée : à ce stade, $\theta$ est une constante déterministe (inconnue) et $I(x)$ un intervalle fixé. C'est avant l'observation que l'intervalle est aléatoire.
\end{remarque}

\begin{remarque}[Confiance vs précision]
Le risque $\alpha$ est typiquement petit : $\alpha = 0{,}05$ (niveau $95\%$) ou $\alpha = 0{,}01$ (niveau $99\%$). Plus $\alpha$ est petit, plus le niveau de confiance est élevé, mais plus l'intervalle est large : c'est un arbitrage inévitable entre \textit{confiance} et \textit{précision}.
\end{remarque}

\begin{remarque}[Bilatéral vs unilatéral]
L'intervalle $[T_-, T_+]$ est dit \textit{bilatéral}. On rencontre aussi des intervalles \textit{unilatéraux}, de la forme $]-\infty,\, T_+(X)]$ (borne de confiance supérieure) ou $[\,T_-(X),\, +\infty[$ (borne de confiance inférieure), utiles lorsqu'on ne s'intéresse qu'à une majoration ou une minoration de $\theta$ (§4.3).
\end{remarque}

\begin{remarque}[Exact vs asymptotique] 
Lorsque la loi exacte des bornes $T_\pm$ est connue sous $P_\theta$, on parle d'IC \textit{exact}. Sinon, on construit des IC \textit{asymptotiques}, pour lesquels la garantie de niveau n'est valable qu'à la limite :

$$
P_\theta\big(\theta \in I_n(X)\big) \xrightarrow[n\to+\infty]{} 1-\alpha.
$$
\end{remarque}

\section{Quantiles}

Toute la construction des intervalles de confiance repose sur les \textbf{quantiles} des lois de probabilité usuelles. On les définit d'abord de façon générale.

\medskip

\begin{mdframed}
\begin{definition}[Fonction quantile] 
Soit $Z$ une variable aléatoire réelle de fonction de répartition $F(z) = P(Z \le z)$. Pour $p \in\, ]0,1[$, le \textit{quantile d'ordre $p$} de $Z$ est

$$
q_p = \inf\{z \in \mathbb R : F(z) \ge p\}.
$$

Lorsque $F$ est continue et strictement croissante (cas des lois à densité qui nous intéressent), $q_p$ est l'unique solution de $F(q_p) = p$, soit $q_p = F^{-1}(p)$.
\end{definition}
\end{mdframed}

\textit{Interprétation : } Le quantile $q_p$ est le seuil en-dessous duquel $Z$ tombe avec probabilité $p$ :

$$
P(Z \le q_p) = p, \qquad P(Z > q_p) = 1-p.
$$

\medskip

\begin{mdframed}
\begin{proposition}[Encadrement par quantiles] 
Pour une loi continue et $\alpha \in\, ]0,1[$,

$$
P\big(q_{\alpha/2} \le Z \le q_{1-\alpha/2}\big) = 1-\alpha,
$$

c'est-à-dire qu'on laisse une masse $\alpha/2$ dans chaque queue. C'est ce découpage \textit{symétrique en probabilité} qui sert par défaut à fabriquer un IC bilatéral.
\end{proposition}
\end{mdframed}

\begin{proof}
$P(q_{\alpha/2} \le Z \le q_{1-\alpha/2}) = F(q_{1-\alpha/2}) - F(q_{\alpha/2}) = (1-\alpha/2) - \alpha/2 = 1-\alpha.$ $\square$
\end{proof}

\medskip

\begin{mdframed}
\begin{notation}[Quantiles des lois usuelles]
On adopte les notations suivantes pour le quantile d'ordre $p$ :

\begin{itemize}
\item $z_p$ :  Normale centrée réduite $\mathcal N(0,1)$
\item $t_{k,p}$ :  Student à $k$ degrés de liberté $\mathcal T(k)$
\item $\chi^2_{k,p}$ :  Khi-deux à $k$ degrés de liberté $\chi^2(k)$
\end{itemize}

Ces valeurs se lisent dans des tables ou se calculent numériquement.
\end{notation}
\end{mdframed}

\begin{remarque}[Lois symétriques] 
Les lois $\mathcal N(0,1)$ et $\mathcal T(k)$ sont \textbf{symétriques} par rapport à $0$, donc leurs quantiles vérifient

$$
q_{1-p} = -\,q_p.
$$

En particulier $z_{1-\alpha/2} = -z_{\alpha/2}$ et $t_{k,1-\alpha/2} = -t_{k,\alpha/2}$. C'est ce qui permet d'écrire l'encadrement bilatéral sous la forme $\pm z_{\alpha/2}$ (resp. $\pm t_{k,\alpha/2}$), où par convention $z_{\alpha/2}$ désigne le quantile d'ordre $1-\alpha/2$, positif. Pour $\alpha = 0{,}05$ : $z_{\alpha/2} = z_{0{,}975} \approx 1{,}96$.
\end{remarque}

\begin{remarque}[Loi du $\chi^2$] 
La loi $\chi^2(k)$ n'est \textbf{pas symétrique} (elle est portée par $\mathbb R_+$). On ne peut donc pas résumer ses deux quantiles par un seul $\pm$ : il faut conserver $\chi^2_{k,\alpha/2}$ et $\chi^2_{k,1-\alpha/2}$ séparément. C'est la source de l'asymétrie des IC pour la variance (§4.6).
\end{remarque}

\section{Construction par pivot}

\subsection{Méthode générale}

\medskip

\begin{mdframed}
\begin{definition}[Pivot]
Une statistique $Z(X,\theta)$ est un \textit{pivot} (ou \textit{quantité pivotale}) si sa loi sous $P_\theta$ ne dépend pas de $\theta$.

\textbf{Méthode :}

\begin{enumerate}
\item \textbf{Pivot :} identifier $Z(X,\theta)$ de loi connue (par exemple $\mathcal N(0,1)$, $\mathcal T(k)$ ou $\chi^2_k$).
\item \textbf{Encadrement :} à l'aide des quantiles, écrire $P\big(q_{\alpha/2} \le Z(X,\theta) \le q_{1-\alpha/2}\big) = 1-\alpha$.
\item \textbf{Inversion :} résoudre l'encadrement en $\theta$ pour obtenir $T_-(X) \le \theta \le T_+(X)$.
\end{enumerate}
\end{definition}
\end{mdframed}

\subsection{Cas gaussien, variance connue}

\begin{exemple}
Soit $(X_1,\dots,X_n)$ i.i.d. de loi $\mathcal N(\theta,\sigma^2)$ avec $\sigma^2$ connue.

\textit{Pivot.} Comme $\overline X_n \sim \mathcal N\!\big(\theta,\sigma^2/n\big)$,

$$
Z = \frac{\sqrt n\,(\overline X_n - \theta)}{\sigma} \sim \mathcal N(0,1),
$$

loi indépendante de $\theta$ : c'est un pivot.

\textit{Encadrement et inversion.}

$$
P_\theta\!\left(-z_{\alpha/2} \le \frac{\sqrt n\,(\overline X_n - \theta)}{\sigma} \le z_{\alpha/2}\right) = 1-\alpha,
$$

d'où l'\textbf{IC exact bilatéral de niveau $1-\alpha$} :

$$
\boxed{\ I(X) = \left[\,\overline X_n - z_{\alpha/2}\,\frac{\sigma}{\sqrt n}\,,\ \overline X_n + z_{\alpha/2}\,\frac{\sigma}{\sqrt n}\,\right].\ }
$$
\end{exemple}

\begin{remarque}
La demi-longueur $z_{\alpha/2}\,\sigma/\sqrt n$ décroît en $1/\sqrt n$ : diviser la largeur par $2$ exige de multiplier $n$ par $4$.
\end{remarque}

\subsection{Cas gaussien, variance inconnue}

Quand $\sigma^2$ est inconnue, le pivot précédent tombe (il contient $\sigma$). On remplace $\sigma$ par un estimateur.

\medskip

\begin{mdframed}
\begin{definition}[Variance empirique corrigée]
$$
S_n^2 = \frac{1}{n-1}\sum_{i=1}^n (X_i - \overline X_n)^2.
$$

Le dénominateur $n-1$ (et non $n$) assure que $S_n^2$ est **sans biais** : $\mathbb E_\theta[S_n^2] = \sigma^2$.
\end{definition}
\end{mdframed}

\medskip

\begin{mdframed}
\begin{proposition}[Loi de Student]
Pour $(X_1,\dots,X_n)$ i.i.d. de loi $\mathcal N(\theta,\sigma^2)$, $\sigma^2$ inconnue,

$$
T_n = \frac{\sqrt n\,(\overline X_n - \theta)}{S_n} \sim \mathcal T(n-1).
$$

La loi de $T_n$ ne dépend pas de $(\theta,\sigma^2)$ : c'est un pivot.
\end{proposition}
\end{mdframed}

\begin{exemple}[IC de Student]
Par symétrie de $\mathcal T(n-1)$, on obtient l'\textbf{IC exact bilatéral} :

$$
\boxed{\ I(X) = \left[\,\overline X_n - t_{n-1,\alpha/2}\,\frac{S_n}{\sqrt n}\,,\ \overline X_n + t_{n-1,\alpha/2}\,\frac{S_n}{\sqrt n}\,\right].\ }
$$
\end{exemple}

\begin{remarque}
Comme $t_{n-1,\alpha/2} > z_{\alpha/2}$, l'IC de Student est toujours plus large que celui à variance connue : c'est le prix de l'ignorance de $\sigma^2$. L'écart s'estompe quand $n \to +\infty$, car $\mathcal T(n-1) \to \mathcal N(0,1)$.
\end{remarque}

\subsection{Intervalles unilatéraux}

Lorsqu'on ne cherche qu'une \textbf{majoration} ou une \textbf{minoration} de $\theta$, on rejette toute la masse $\alpha$ dans une seule queue plutôt que $\alpha/2$ dans chacune. On utilise alors le quantile d'ordre $1-\alpha$ (et non $1-\alpha/2$).

\begin{exemple}[Bornes de confiance, cas gaussien à variance connue] 
Partant du même pivot $Z = \sqrt n\,(\overline X_n - \theta)/\sigma \sim \mathcal N(0,1)$ :

\begin{itemize}
\item \textbf{Borne de confiance supérieure} : $P_\theta(Z \ge -z_\alpha) = 1-\alpha$ donne

$$
P_\theta\!\left(\theta \le \overline X_n + z_\alpha\,\frac{\sigma}{\sqrt n}\right) = 1-\alpha,
\qquad
I(X) = \left]-\infty,\ \overline X_n + z_\alpha\,\frac{\sigma}{\sqrt n}\right].
$$
\item \textbf{Borne de confiance inférieure} :

$$
I(X) = \left[\,\overline X_n - z_\alpha\,\frac{\sigma}{\sqrt n},\ +\infty\right[.
$$
\end{itemize}

Ici $z_\alpha$ est le quantile d'ordre $1-\alpha$ (pour $\alpha = 0{,}05$ : $z_\alpha \approx 1{,}645$, à comparer à $z_{\alpha/2} \approx 1{,}96$).
\end{exemple}

\begin{remarque} 
Une borne unilatérale est \textbf{plus resserrée} du côté considéré que l'IC bilatéral, puisque $z_\alpha < z_{\alpha/2}$. On l'emploie quand la question est directionnelle (« $\theta$ dépasse-t-il un seuil ? »). Le même principe s'applique avec les quantiles de Student ou du $\chi^2$.
\end{remarque}

\section{Intervalles de confiance asymptotiques}

\subsection{Construction par normalité asymptotique}

Quand la loi exacte de l'estimateur est inconnue, on exploite la normalité asymptotique et les quantiles de $\mathcal N(0,1)$.

\medskip

\begin{mdframed}
\begin{proposition}[IC asymptotique] 
Soit $\hat\theta_n$ asymptotiquement normal :

$$
\sqrt n\,(\hat\theta_n - \theta) \xrightarrow[n\to+\infty]{\text{loi}} \mathcal N\big(0,\sigma^2(\theta)\big),
$$

et $\hat\sigma_n^2$ un estimateur consistant de $\sigma^2(\theta)$. Alors, pour tout $\theta \in \Theta$,

$$
P_\theta\!\left(\theta \in \left[\hat\theta_n - z_{\alpha/2}\,\frac{\hat\sigma_n}{\sqrt n},\ \hat\theta_n + z_{\alpha/2}\,\frac{\hat\sigma_n}{\sqrt n}\right]\right) \xrightarrow[n\to+\infty]{} 1-\alpha.
$$
\end{proposition}
\end{mdframed}

\begin{proof}
Normalité asymptotique + consistance de $\hat\sigma_n$ + lemme de Slutsky donnent $Z_n = \sqrt n\,(\hat\theta_n - \theta)/\hat\sigma_n \to \mathcal N(0,1)$ en loi. Donc $P_\theta(|Z_n| \le z_{\alpha/2}) \to P(|Z| \le z_{\alpha/2}) = 1-\alpha$, et on isole $\theta$. $\square$
\end{proof}

\begin{remarque}
La qualité de l'approximation dépend de la vitesse de convergence de $Z_n$ vers $\mathcal N(0,1)$. Pour $n$ petit ou une loi très asymétrique, elle peut être médiocre ; le *bootstrap* permet alors d'améliorer la précision.
\end{remarque}

\begin{exemple}[Proportion dans un modèle de Bernoulli]
Soit $(X_1,\dots,X_n)$ i.i.d. de loi $\mathcal B(\theta)$, $\theta \in\, ]0,1[$, et $\hat\theta_n = \overline X_n$. Par le TCL,

$$
\sqrt n\,(\hat\theta_n - \theta) \xrightarrow[n\to+\infty]{\text{loi}} \mathcal N\big(0,\ \theta(1-\theta)\big).
$$

On remplace la variance asymptotique $\sigma^2(\theta) = \theta(1-\theta)$ par le plug-in consistant $\hat\sigma_n^2 = \hat\theta_n(1-\hat\theta_n)$ :

$$
I_n = \left[\,\hat\theta_n - z_{\alpha/2}\sqrt{\frac{\hat\theta_n(1-\hat\theta_n)}{n}}\,,\ \hat\theta_n + z_{\alpha/2}\sqrt{\frac{\hat\theta_n(1-\hat\theta_n)}{n}}\,\right].
$$

Pour $\alpha = 0{,}05$, $z_{\alpha/2} \approx 1{,}96$.
\end{exemple}

\begin{remarque}
Très utilisé (sondages, essais cliniques). La largeur est maximale en $\theta = 1/2$, d'où la borne conservative : demi-largeur $\le z_{\alpha/2}/(2\sqrt n)$. Pour $n = 1000$ cela donne $\approx 0{,}031$, nettement plus fin que la borne grossière $0{,}0707$ de Tchebychev au chapitre 1.
\end{remarque}

\section{Longueur d'un intervalle et borne de Cramér-Rao}

Cette section relie la précision d'un IC à la borne de Cramér-Rao du chapitre 3.

La demi-longueur d'un IC asymptotique est

$$
\ell_n = z_{\alpha/2}\,\sqrt{\frac{\sigma^2(\theta)}{n}}.
$$

Deux leviers agissent sur elle :

\begin{itemize}
\item La taille d'échantillon : $\ell_n$ décroît en $1/\sqrt n$ (précision lente).
\item La variance asymptotique : à $n$ et $\alpha$ fixés, plus $\sigma^2(\theta)$ est petite, plus l'IC est court. 

Minimiser $\sigma^2(\theta)$, c'est produire l'IC le plus précis.
\end{itemize}

\textit{Lien avec Cramér-Rao.} Dans un modèle régulier (hypothèses (H1)–(H4) du chapitre 3), tout estimateur régulier vérifie

$$
\sigma^2(\theta) \ge \frac{1}{I(\theta)},
$$

où $I(\theta)$ est l'information de Fisher d'une observation. Donc

$$
\ell_n \ge \frac{z_{\alpha/2}}{\sqrt{n\,I(\theta)}}.
$$

\medskip

\begin{mdframed}
\begin{definition}[Efficacité asymptotique] 
Un estimateur est \textit{asymptotiquement efficace} s'il atteint cette borne, $\sigma^2(\theta) = 1/I(\theta)$. Son IC asymptotique est alors le **plus court possible** à niveau donné :

$$
I_n = \left[\,\hat\theta_n - \frac{z_{\alpha/2}}{\sqrt{n\,I(\hat\theta_n)}}\,,\ \hat\theta_n + \frac{z_{\alpha/2}}{\sqrt{n\,I(\hat\theta_n)}}\,\right].
$$
\end{definition}
\end{mdframed}

\begin{remarque}
Sous conditions de régularité, l'EMV est asymptotiquement efficace : $\sqrt n\,(\hat\theta_{\text{EMV}} - \theta) \to \mathcal N\big(0, 1/I(\theta)\big)$. C'est la justification profonde de son usage : il fournit asymptotiquement l'IC le plus précis. Une information de Fisher grande (log-vraisemblance « pointue ») donne un IC court, une information faible un IC large.
\end{remarque}

\begin{exemple}[Bernoulli, suite] 
Ici 

$$
I(\theta) = 1/[\theta(1-\theta)],
$$ 

donc 

$$
1/I(\theta) = \theta(1-\theta) = \sigma^2(\theta)
$$ 

La moyenne empirique est efficace et son IC asymptotiquement optimal.
\end{exemple}

\section{Intervalles de confiance dans les modèles gaussiens}

\subsection{Vecteur gaussien}

\medskip

\begin{mdframed}
\begin{definition}
Un vecteur aléatoire $X$ à valeurs dans $\mathbb R^k$ est un \textit{vecteur gaussien} ssi toute combinaison linéaire de ses coordonnées est une variable aléatoire réelle gaussienne : pour tout $c \in \mathbb R^k$, $c^\top X = \sum_{i=1}^k c_i X_i \sim \mathcal N(\mu,\sigma^2)$ pour certains $\mu \in \mathbb R$, $\sigma^2 \in \mathbb R_+$.
\end{definition}
\end{mdframed}

\medskip

\begin{mdframed}
\begin{proposition}
Un vecteur gaussien $X$ est caractérisé par son vecteur des espérances $\mathbb E[X]$ et sa matrice de covariance $V = \mathbb V[X]$, d'élément $(i,j)$ égal à $V_{i,j} = \mathrm{Cov}(X_i,X_j)$. On note $X \sim \mathcal N\big(\mathbb E[X], V\big)$.
\end{proposition}
\end{mdframed}

\begin{remarque}
Un $n$-échantillon i.i.d. de loi $\mathcal N(\mu,\sigma^2)$ est un vecteur gaussien de loi $\mathcal N(\mu\,\mathbf 1_n,\ \sigma^2 I_n)$, où $\mathbf 1_n = (1,\dots,1)^\top \in \mathbb R^n$ : ses $n$ composantes sont indépendantes, d'espérance $\mu$ et de variance $\sigma^2$.
\end{remarque}

\subsection{Théorème de Cochran}

Soit $(X_1,\dots,X_n)$ i.i.d. de loi $\mathcal N(\mu,\sigma^2)$. On pose
$$
\overline X_n = \frac{1}{n}\sum_{i=1}^n X_i,
\qquad
S_n^2 = \frac{1}{n-1}\sum_{i=1}^n (X_i - \overline X_n)^2.
$$

\medskip

\begin{mdframed}
\begin{theoreme}[Cochran]
\begin{enumerate}
\item $\overline X_n \sim \mathcal N\!\big(\mu,\ \sigma^2/n\big)$ ;
\item $\dfrac{(n-1)S_n^2}{\sigma^2} \sim \chi^2(n-1)$ ;
\item $\overline X_n$ et $S_n^2$ sont \textit{indépendants}.
\end{enumerate}
\end{theoreme}
\end{mdframed}

\begin{proof}[Idée]
$X = (X_1,\dots,X_n)$ est gaussien, de loi $\mathcal N(\mu\mathbf 1_n,\sigma^2 I_n)$. Soit $D$ la droite dirigée par $\mathbf 1_n$ et $H = D^\perp$ son orthogonal. La projection de $X$ sur $D$ fait apparaître $\overline X_n$, celle sur $H$ les résidus $(X_i - \overline X_n)_i$. Deux projections d'un vecteur gaussien sur des sous-espaces orthogonaux sont indépendantes et gaussiennes (d'où le point 3). La norme au carré de la projection sur $H$ vaut $\sum_i (X_i - \overline X_n)^2/\sigma^2 = (n-1)S_n^2/\sigma^2$, somme des carrés de $\dim H = n-1$ variables $\mathcal N(0,1)$ indépendantes : c'est une loi $\chi^2(n-1)$ (point 2). Le point 1 est immédiat par linéarité. $\square$
\end{proof}

\medskip

\begin{mdframed}
\begin{corollaire}[Pivot de Student]
En combinant les trois points et en normalisant par $S_n$ au lieu de $\sigma$ :

$$
\frac{\sqrt n\,(\overline X_n - \mu)}{S_n} \sim \mathcal T(n-1),
$$
\end{corollaire}
\end{mdframed}

\subsection{Intervalles de confiance pour $\mu$}

\begin{itemize}
\item $\sigma^2$ connue (pivot $\mathcal N(0,1)$, point 1 de Cochran) :

$$
I(X) = \left[\,\overline X_n - z_{\alpha/2}\,\frac{\sigma}{\sqrt n}\,,\ \overline X_n + z_{\alpha/2}\,\frac{\sigma}{\sqrt n}\,\right].
$$

\item $\sigma^2$ inconnue (pivot de Student) :

$$
I(X) = \left[\,\overline X_n - t_{n-1,\alpha/2}\,\frac{S_n}{\sqrt n}\,,\ \overline X_n + t_{n-1,\alpha/2}\,\frac{S_n}{\sqrt n}\,\right].
$$
\end{itemize}

\subsection{Intervalles de confiance pour $\sigma^2$}

Le point 2 de Cochran fournit un pivot pour $\sigma^2$ :

$$
Z = \frac{(n-1)S_n^2}{\sigma^2} \sim \chi^2(n-1),
$$

de loi indépendante de $(\mu,\sigma^2)$. La loi du $\chi^2$ étant \textbf{asymétrique}, on prend deux quantiles distincts :

$$
P_\theta\!\left(\chi^2_{n-1,\,\alpha/2} \le \frac{(n-1)S_n^2}{\sigma^2} \le \chi^2_{n-1,\,1-\alpha/2}\right) = 1-\alpha.
$$

On isole $\sigma^2$ (l'inversion \textbf{renverse} les inégalités, $\sigma^2$ étant au dénominateur) :

$$
\boxed{\ I(X) = \left[\ \frac{(n-1)S_n^2}{\chi^2_{n-1,\,1-\alpha/2}}\,,\ \frac{(n-1)S_n^2}{\chi^2_{n-1,\,\alpha/2}}\ \right].\ }
$$

En prenant les racines des bornes (positives), on obtient l'IC pour l'écart-type $\sigma$ :

$$
I(X) = \left[\ \sqrt{\frac{(n-1)S_n^2}{\chi^2_{n-1,\,1-\alpha/2}}}\,,\ \sqrt{\frac{(n-1)S_n^2}{\chi^2_{n-1,\,\alpha/2}}}\ \right].
$$

\begin{remarque}
Contrairement au cas de $\mu$, cet intervalle n'est **pas symétrique** autour de $S_n^2$, du fait de l'asymétrie de la loi du $\chi^2$. Le découpage symétrique en probabilité ($\alpha/2$ dans chaque queue) est le choix standard, mais pas celui qui minimise strictement la longueur.
\end{remarque}
 

\chapter*{TD 7 - Intervalles de confiance}

\setcounter{exercice}{0}

\section*{Rappels}
\begin{enumerate}
    \item On dit qu'une variable aléatoire $U$ suit la loi du $\chi^2$ à $d$ degrés de liberté, si
    \begin{equation*}
    U \overset{\text{(loi)}}{=} Z_1^2 + \dots + Z_d^2
    \end{equation*}
    où les $Z_i$ sont i.i.d de loi $\mathcal{N}(0, 1)$. On note dans ce cas $U \sim \chi^2(d)$.\\
    On peut montrer (voir l'exercice ci-dessous) que si $U \sim \chi^2(d)$, alors $U$ admet pour densité
    \begin{equation*}
    f_U(x) = C_d \exp(-x/2) x^{d/2 - 1}, \quad \forall x \ge 0.
    \end{equation*}
    
    \item On dit qu'une variable aléatoire $T$ suit la loi de Student à $d$ degrés de liberté si
    \begin{equation*}
    T \overset{\text{(loi)}}{=} \frac{Z}{\sqrt{U/d}}
    \end{equation*}
    où $Z \sim \mathcal{N}(0, 1)$ et $U \sim \chi^2(d)$ sont indépendantes. On peut montrer que dans ce cas $T$ admet pour densité
    \begin{equation*}
    f_T(x) = C'_d \left( 1 + \frac{x^2}{d} \right)^{-(d+1)/2}, \quad \forall x \in \mathbb{R}.
    \end{equation*}
\end{enumerate}

\smallskip

\begin{mdframed}
\begin{exercice}
Soit $X$ une variable aléatoire à valeurs dans $\mathbb{N}$ définie comme l'instant de premier succès dans un schéma de Bernoulli de paramètre $q \in \, ]0, 1[$. On dispose de $X_1, \dots, X_n$ un échantillon indépendant de taille $n$ de même loi que $X$.

\begin{enumerate}
    \item Donner la loi de $X$.
    \item Déterminer $\hat{q}_n$, l'estimateur du maximum de vraisemblance de $q$.
    \item Montrer que l'estimateur du maximum de vraisemblance est asymptotiquement normal.
    \item Donner un intervalle de confiance pour $q$ de niveau $1 - \alpha$.
    \item Une société de transport en commun par bus veut estimer le nombre de passagers ne validant pas leur titre de transport sur une ligne de bus déterminée. Elle dispose pour cela, pour un jour de semaine moyen, du nombre $n_0$ de tickets compostés sur la ligne et des résultats de l'enquête suivante : à chacun des arrêts de bus de la ligne, des contrôleurs comptent le nombre de passagers sortant des bus et ayant validé leur ticket jusqu'à la sortie du premier fraudeur. On a les données (simulées) suivantes :
    \begin{center}
    \begin{tabular}{cccccccccc}
    44 & 07 & 26 & 34 & 09 & 21 & 219 & 05 & 11 & 90 \\
    02 & 07 & 59 & 38 & 57 & 15 & 81  & 01 & 11 & 06 \\
    44 & 15 & 24 & 129& 19 & 22 & 69  & 14 & 89 & 29 \\
    12 & 18 & 10 & 19 & 21 & 02 & 24  & 37 & 28 & 156
    \end{tabular}
    \end{center}
    \item Estimer la probabilité de fraude. Donner un intervalle de confiance de niveau 0.95\%. Estimer le nombre de fraudeurs si $n_0 = 20000$.
\end{enumerate}
\end{exercice}
\end{mdframed}

\smallskip

\begin{mdframed}
\begin{exercice}
Soit $X_1, \dots, X_n$ une suite de variables aléatoires indépendantes et identiquement distribuées de loi de Bernoulli $p \in [0, 1]$. On pose $\hat{p}_n = \frac{1}{n} \sum_{i=1}^n X_i$.

\begin{enumerate}
    \item Montrer l'inégalité $\text{Var}(\hat{p}_n) \le \frac{1}{4n}$.
    \item Un institut de sondage souhaite estimer avec une précision de 3 points (à droite et à gauche) la probabilité qu'un individu vote pour le maire actuel aux prochaines élections. Combien de personnes est-il nécessaire de sonder ?
    \item Sur un échantillon représentatif de 1000 personnes, on rapporte les avis favorables pour un homme politique. En novembre, il y avait 38\% d'avis favorables, en décembre 36\%. Un éditorialiste dans son journal prend très au sérieux cette chute de 2 points d'un futur candidat. Confirmer ou infirmer la position du journaliste.
    \item Reprendre la question 2, et établir un intervalle de confiance à l'aide de l'inégalité de Hoeffding : si $X_1, \dots, X_n$ sont des variables aléatoires i.i.d à valeurs dans $[a, b]$, alors
    \begin{equation}
    \mathbb{P}(|\bar{X}_n - \mathbb{E}[X_1]| > t) \le 2 \exp \left( -\frac{2nt^2}{(b - a)^2} \right), \quad \forall t > 0.
    \end{equation}
\end{enumerate}
\end{exercice}
\end{mdframed}

\smallskip

\begin{mdframed}
\begin{exercice}
\begin{enumerate}
    \item Soit $U$ une variable aléatoire suivant la loi du $\chi^2$ à $d$ degrés de libertés. Donner l'espérance de $U$. Montrer que $U$ admet la densité suivante
    \begin{equation*}
    f_U(y) = C_d \exp(-y/2) y^{d/2 - 1}, \quad \forall y \ge 0,
    \end{equation*}
    où $C_d$ est une constante de renormalisation.\\
    \textit{Indication : utiliser la propriété de stabilité des lois Gamma.}
    
    \item Soit $X_1, \dots, X_n$ des variables aléatoires gaussiennes i.i.d de loi $\mathcal{N}(m, \sigma^2)$.\\
    Le but de cette question est de montrer que les variables $\bar{X}$ et $V$ définies par
    \begin{equation*}
    \bar{X} = \frac{1}{n}\sum_{i=1}^n X_i \quad \text{et} \quad V = \frac{1}{n}\sum_{i=1}^n (X_i - \bar{X})^2
    \end{equation*}
    sont indépendantes et que $\frac{n}{\sigma^2}V$ suit une loi du $\chi^2$ à $n - 1$ degrés de liberté.
    \begin{enumerate}
        \item[i--] Montrer qu'on peut se ramener au cas où $m = 0$ et $\sigma^2 = 1$.\\
        On suppose désormais que $m = 0$ et $\sigma^2 = 1$.
        \item[ii--] Montrer que
        \begin{equation*}
        V = \frac{1}{n} \left( \sum_{i=1}^n X_i^2 - \left( \frac{1}{\sqrt{n}} \sum_{i=1}^n X_i \right)^2 \right).
        \end{equation*}
        \item[iii--] On se donne une matrice $M$ de taille $n \times n$ orthogonale (i.e. $M^t \cdot M = I$) telle que la première ligne de $M$ soit $[\frac{1}{\sqrt{n}}, \dots, \frac{1}{\sqrt{n}}]$. On pose $Y = MX$, avec $X = [X_1, \dots, X_n]^t$. Montrer que $Y$ est un vecteur gaussien standard et que $\sum_{i=1}^n Y_i^2 = \sum_{i=1}^n X_i^2$.
        \item[iv--] Conclure.
    \end{enumerate}
\end{enumerate}
\end{exercice}
\end{mdframed}

\smallskip

\begin{mdframed}
\begin{exercice}
Soit $X_1, \dots, X_n$ un échantillon de loi $\mathcal{N}(m, \sigma^2)$ avec $m$ et $\sigma^2$ inconnus. On pose
\begin{equation*}
\bar{X} = \frac{1}{n}\sum_{i=1}^n X_i \quad \text{et} \quad \hat{\sigma}^2 = \frac{1}{n-1}\sum_{i=1}^n (X_i - \bar{X})^2
\end{equation*}

\begin{enumerate}
    \item Montrer que la variable $T$ définie par
    \begin{equation*}
    T = \sqrt{n} \frac{\bar{X} - m}{\hat{\sigma}}
    \end{equation*}
    suit une loi de Student à $n - 1$ degrés de liberté.
    \item Montrer que si $t_\alpha$ est tel que $\mathbb{P}(|T| \le t_\alpha) = 1 - \alpha$, alors
    \begin{equation*}
    \mathbb{P}\left( m \in \left[ \bar{X} - \frac{\hat{\sigma}t_\alpha}{\sqrt{n}}, \bar{X} + \frac{\hat{\sigma}t_\alpha}{\sqrt{n}} \right] \right) = 1 - \alpha.
    \end{equation*}
\end{enumerate}
\end{exercice}
\end{mdframed}

\smallskip

\begin{mdframed}
\begin{exercice}
On désire estimer la production d'une nouvelle espèce de pommier. On suppose que la production d'un pommier de cette espèce suit une loi normale d'espérance $\mu$ et d'écart-type $\sigma$ inconnus.

\begin{enumerate}
    \item Rappeler les estimateurs de $\mu$ et de $\sigma^2$ dans ce modèle, et indiquer les méthodes qui y conduisent.
    \item Sur un échantillon de 15 pommiers, on a observé une récolte moyenne de 52 kg avec un écart-type empirique de 5 kg. Donner un intervalle de confiance pour la production moyenne des pommiers de cette espèce, de niveau 0.95 puis 0.99.
    \item Donner un intervalle de confiance pour l'écart-type $\sigma$, de niveau 0.95.
    \item Sur un échantillon de 80 pommiers, on observe une récolte moyenne de 51.5 kg avec un écart-type de 4.5 kg. Donner un intervalle de confiance pour la production de cette espèce, de niveau 0.95 puis 0.99.
\end{enumerate}
\end{exercice}
\end{mdframed}
\chapter*{TD 8 - Exercices Complet}

\setcounter{exercice}{0}

\begin{mdframed}
\begin{exercice}
Soit $X_1, \dots, X_n$ un échantillon i.i.d. de loi exponentielle de densité donnée pour tout $x \in \mathbb{R}$ par
\begin{equation*}
f_\lambda(x) = \lambda e^{-\lambda x} \mathbf{1}_{[0, +\infty[}(x),
\end{equation*}
de paramètre $\lambda > 0$ inconnu.

\begin{enumerate}
    \item Calculer la fonction de répartition de $X_1$.
    \item On note $m$ la médiane de $X_1$. Montrer que $m = \ln(2)/\lambda$.
    \item On note $\hat{m}_n$ la médiane empirique de l'échantillon. Rappeler la définition de $\hat{m}_n$ et donner la loi limite de $\sqrt{n}(\hat{m}_n - m)$ lorsque $n \to \infty$.
    \item On pose $\hat{\lambda}_n = \ln(2)/\hat{m}_n$. Donner la loi limite de $\sqrt{n}(\hat{\lambda}_n - \lambda)$.
\end{enumerate}
\end{exercice}
\end{mdframed}

\smallskip

\begin{mdframed}
\begin{exercice}
Soit $0 < a < b$ et $\theta \in [0, 1]$, on considère un $n$-échantillon $(X_1, \dots, X_n)$ de loi ayant pour densité $p_\theta$ définie pour tout $x \in \mathbb{R}$, par
\begin{equation*}
p_\theta(x) = \theta \frac{1}{a}\mathbf{1}_{[0, a]}(x) + (1 - \theta) \frac{1}{b}\mathbf{1}_{[0, b]}(x).
\end{equation*}
Dans la suite, on suppose que les paramètres $a$ et $b$ sont connus et on cherche à estimer le paramètre $\theta$ inconnu.

\begin{enumerate}
    \item Montrer que la vraisemblance est donnée par
    \begin{equation*}
    L(X_1, \dots, X_n; \theta) = \left( \frac{\theta}{a} + \frac{1 - \theta}{b} \right)^N \left( \frac{1 - \theta}{b} \right)^{n - N},
    \end{equation*}
    où $N = \text{Card}\{i \in \{1, \dots, n\} : X_i \in [0, a]\}$.
    \item Déterminer l'estimateur $\hat{\theta}$ du maximum de vraisemblance et calculer son biais.
    \item Calculer $\mathbb{E}[X_1]$ et en déduire un estimateur $\tilde{\theta}$ construit à l'aide de la méthode des moments et calculer son biais.
    \item Calculer le risque quadratique des estimateurs $\hat{\theta}$ et $\tilde{\theta}$. Comparer ces risques quadratiques quand $b = 2a$.
\end{enumerate}
\end{exercice}
\end{mdframed}

\smallskip

\begin{mdframed}
\begin{exercice}
Pour tous $a \in \mathbb{R}$ et $b > 0$, on note $\mathcal{P}(a, b)$ la loi sur $\mathbb{R}$ ayant pour densité
\begin{equation*}
p_{a,b}(x) = \frac{b}{(x - a)^{b+1}} \mathbf{1}_{[a+1, +\infty[}(x).
\end{equation*}

\subsection*{Partie I : Préliminaires}
\begin{enumerate}
    \item Vérifier que $p_{a,b}$ est bien une densité de probabilité sur $\mathbb{R}$.
    \item Calculer la fonction de répartition $F_{a,b}$ de $\mathcal{P}(a, b)$.
    \item Montrer que si $X \sim \mathcal{P}(a, b)$, alors $\ln(X - a)$ suit une loi exponentielle de paramètre $b$.
\end{enumerate}

\subsection*{Partie II : Estimation de $b$ dans le cas où $a$ est connu}
Dans cette partie, on considère un $n$-échantillon $(X_1, \dots, X_n)$ de loi $\mathcal{P}(a, b)$ et on suppose que $a$ est connu et $b$ inconnu.
\begin{enumerate}
    \setcounter{enumi}{3}
    \item Montrer que l'estimateur du maximum de vraisemblance de $b$ est donné par
    \begin{equation*}
    \hat{b}_n = \frac{n}{\sum_{i=1}^n \ln(X_i - a)}.
    \end{equation*}
    \item Justifier que $\hat{b}_n$ est consistant.
    \item Donner la limite en loi de $\sqrt{n}(\hat{b}_n - b)$ lorsque $n \to + \infty$.
\end{enumerate}

\subsection*{Partie III : Estimation de $a$ dans le cas où $b$ est connu}
Dans cette partie, on suppose que $a$ est inconnu et $b$ est connu et on cherche à estimer le paramètre $a$.
\begin{enumerate}
    \setcounter{enumi}{6}
    \item Montrer que l'estimateur du maximum de vraisemblance de $a$ est donné par
    \begin{equation*}
    \hat{a}_n = \min(X_1, \dots, X_n) - 1.
    \end{equation*}
    \item Calculer la fonction de répartition de $\hat{a}_n - a$. En déduire que $\hat{a}_n$ est consistant.
    \item En utilisant la question précédente, donner un intervalle de confiance non asymptotique de niveau $1 - \alpha$, $\alpha \in \, ]0, 1[$, pour estimer $a$.
\end{enumerate}
\end{exercice}
\end{mdframed}

\smallskip

\begin{mdframed}
\begin{exercice}
Soient $X_1, \dots, X_n$, des variables aléatoires i.i.d. de densité donnée pour tout $x \in \mathbb{R}$ par
\begin{equation*}
f(x) = \frac{\theta}{x^{\theta+1}} \mathbf{1}_{\{x \ge 1\}},
\end{equation*}
où $\theta$ est un paramètre inconnu.

\begin{enumerate}
    \item Dans cette question, on suppose d'abord que l'ensemble des paramètres est $\Theta = \{\theta > 1\}$. Calculer $\mathbb{E}_\theta[X_1]$ puis estimer $\theta$ par la méthode des moments.
    \item Dans cette question, on suppose désormais que l'ensemble des paramètres est $\Theta = \{\theta > 0\}$.
    \begin{enumerate}
        \item[i--] Montrer qu'il n'existe pas de $r > 0$ tel que la fonction $\theta \in \mathbb{R}_+^* \mapsto \mathbb{E}_\theta[X_1^r]$ soit injective.
        \item[ii--] Montrer que $\mathbb{E}_\theta\left[ \frac{1}{X_1} \right] = \frac{\theta}{\theta+1}$ puis estimer $\theta$ par la méthode des moments. On notera $\hat{\theta}_n^{MM}$ l'estimateur obtenu.
        \item[iii--] $\hat{\theta}_n^{MM}$ est-il consistant ?
        \item[iv--] Donner la loi asymptotique de $\sqrt{n}(\hat{\theta}_n^{MM} - \theta)$ lorsque $n \to +\infty$.
        \item[v--] Montrer que l'estimateur du maximum de vraisemblance $\hat{\theta}_n^{MV}$ vaut
        \begin{equation*}
        \hat{\theta}_n^{MV} = \frac{n}{\sum_{i=1}^n \ln(X_i)}.
        \end{equation*}
        \item[vi--] $\hat{\theta}_n^{MV}$ est-il consistant ?\\
        \textit{Indication : on pourra utiliser le changement de variables $u = \ln(x)$.}
    \end{enumerate}
\end{enumerate}
\end{exercice}
\end{mdframed}

\smallskip

\begin{mdframed}
\begin{exercice}
Soient $X_1, \dots, X_n$ des variables aléatoires i.i.d. de densité donnée pour tout $x \in \mathbb{R}$ par
\begin{equation*}
f(x) = \frac{1}{2\sqrt{x\theta}} \mathbf{1}_{]0, \theta]}(x),
\end{equation*}
où $\theta > 0$ est un paramètre inconnu.

\begin{enumerate}
    \item Calculer $\mathbb{E}_\theta[X_1]$.
    \item En déduire un estimateur $\hat{\theta}_n$ par la méthode des moments.
    \item Déterminer l'estimateur $\tilde{\theta}_n$ de $\theta$ par maximum de vraisemblance.
    \item Déterminer la loi de $\tilde{\theta}_n$.
    \item Montrer que le risque quadratique de $\tilde{\theta}_n$ tend vers 0 quand $n$ tend vers $+\infty$.
    \item Comparer les risques quadratiques de $\hat{\theta}_n$ et de $\tilde{\theta}_n$.
    \item Préciser la loi asymptotique de $\sqrt{n}(\hat{\theta}_n - \theta)$ et de $n(\theta - \tilde{\theta}_n)$.
    \item Montrer que la médiane de $X_1$ est $F_\theta^{-1}(1/2) = \theta/4$. En déduire un estimateur $\hat{\theta}_n^m$ consistant pour l'estimation de $\theta$.
    \item Donner la loi limite de $\hat{\theta}_n^m - \theta$ après une renormalisation convenable.
\end{enumerate}
\end{exercice}
\end{mdframed}

\smallskip

\begin{mdframed}
\begin{exercice}
Soient $X_1, \dots, X_n$ des variables aléatoires i.i.d. de densité donnée pour tout $x \in \mathbb{R}$ par
\begin{equation*}
f(x) = e^{-\theta} \mathbf{1}_{]0, e^\theta]}(x),
\end{equation*}
où $\theta > 0$ est un paramètre inconnu.

\begin{enumerate}
    \item Déterminer l'estimateur du maximum de vraisemblance de $\theta$.
    \item Montrer que $\theta - \ln(X_{(n)})$ suit une loi exponentielle de paramètre $n$ (où $X_{(n)} = \max(X_1, \dots, X_n)$).
    \item Soit $0 < \alpha < 1$. Montrer que 
    \begin{equation*}
    \left[ \ln(X_{(n)}) \, , \, \ln(X_{(n)}) - \frac{1}{n} \ln(\alpha) \right]
    \end{equation*}
    est un intervalle de confiance pour $\theta$ de niveau $1 - \alpha$.
\end{enumerate}
\end{exercice}
\end{mdframed}

\smallskip

\begin{mdframed}
\begin{exercice}

\textbf{Rappels}

\smallskip

\begin{itemize}
\item Une v.a.  de loi exponentielle $A \sim \mathcal{E}(\theta)$ , $\theta >0$ admet pour densité :

$$
f_A(x) = \theta e^{-\theta x} \mathbb{1}_{x>0}
$$
\item Soit $A_1,\cdots,A_n$ un $n$-échantillon i.i.d. de loi $\mathcal{E}(\theta)$. Alors $S = \sum_{i=1}^n A_i$ suit une loi Gamma $\Gamma(n,\theta)$ de densité donnée par

$$
f_S(x) = \frac{\theta^n}{\Gamma(n)} x^{n-1} e^{- \theta x} \mathbb{1}_{x>0}
$$

où $\Gamma$ désigne la fonction Gamma. 
\item Pour tout entier entier $n \in \mathbb{N}^*$, $\Gamma(n) = (n-1)!$
\end{itemize}

\smallskip

\textbf{Questions préliminaires}

\smallskip

\begin{enumerate}
\item Soit $A \sim \mathcal{E}(\theta)$ avec $\theta > 0$. Soit $A_1, \cdots A_n$, $n$ copies i.i.d. de $A$. Donner la vraisemblance associée aux observation $A_1, \cdots A_n$.
\item Soit $B \sim \mathcal{B}(p)$ avec $p \in ]0,1[$. Soit $B_1, \cdots B_n$, $n$ copies i.i.d. de $B$. Donner la vraisemblance associée aux observation $B_1, \cdots B_n$.
\end{enumerate}

\smallskip

\textbf{Problème}

\smallskip

On s'interesse à la durée de vie de deux composants électroniques se trouvant dans un système solidaire. Si l'un des deux composants tombe en panne, le système tout entier doit être changé.

Les durée de vie des deux composants sont modélisés par des variables aléatoires exponentielles de paramètres $\lambda > 0$ et $\mu > 0$ indépendantes. Pour $n$ systèmes, on a donc un $n$-échatillon i.i.d. $X_1,\cdots,X_n$ de loi $\mathcal{E}(\lambda)$ et un $n$-échantillon i.i.d. $Y_1, \cdots , Y_n$ de loi $\mathcal{E}(\mu)$, indépendant du précédent.

Cépendant, on observe uniquement la durée de vie du commposant tombé en panne, ce qui veut dire que l'on observe uniquement le $n$-échantillon 

$(Z_1,W_1) , \cdots , (Z_n,W_n)$, où pour tout $i \in 1 , \cdots , n$,

$$
Z_i = \min (X_i,Y_i) \qquad W_i = \mathbb{1}_{Z_i = X_i} = \mathbb{1}_{X_i \leq Y_i}   
$$

\smallskip

\begin{enumerate}
\item Calculer la fonction de répartition de $Z_i$ et en déduire $Z_i \sim \mathcal{E} (\lambda + \mu )$
\item Montrer que $W_i \sim \mathcal{B}( \frac{\lambda}{\lambda + \mu} )$
\item
\begin{enumerate}
\item Soit $z \geq 0$. Montrer que $\mathbb{P}(Z_i \leq z , W_i = 1) = \frac{\lambda}{\lambda + \mu} \mathbb{P}(Z_i \leq z_i)$.
\item En conclure que les v.a. $Z_i$ et $W_i$ sont indépendantes.
\end{enumerate}
\end{enumerate}

On rappelle que par indépendance entre les $Z_i$ et les $X_i$, la vraisemblance du $n$-échantillon i.i.d. $(Z_1,W_1) , \cdots , (Z_n,W_n)$ est définie comme le produit de la vraisemblance du $n$-échantillon $Z_1 , \cdots , Z_n$ par la vraisemblance du $n$-échantillon $W_1 , \cdots , W_n$. On suppose désormais que le paramètre $\mu$ \textbf{est connu.}

\begin{enumerate}
\item Déterminer l'EMV $\hat{\lambda}_n$ de $\lambda$.
\item Montrer que $\hat{\lambda}_n$ est consistant.
\item Calculer le biais de $\hat{\lambda}_n$ puis montrer qu'il tend ver 0 lorsque $n \to \infty$.
\item Déterminer la loi asymptotique de $\sqrt{n} ( \hat{\lambda}_n - \lambda)$ lorsque $n \to \infty$.
\item Donner un intervalle de confiance bilatéral symétrique asymptotique de niveau $1 - \alpha$ pour $\lambda$. Dans la question suivantee on suppose que $\lambda$ et $\mu$ \textbf{sont inconnus.}
\item Déterminer l'EMV $(\hat{\lambda}_n,\hat{\mu}_n)$ de $( \lambda , \mu)$.
\end{enumerate}
\end{exercice}
\end{mdframed}

\smallskip

\begin{mdframed}
\begin{exercice}
Soit $X_1, \cdots, X_n$ des variables aléatoires i.i.d. de deinsité donnée, pour tout $x \in \mathbb{R}$, par :

$$
f_\theta(x) = C_\theta x^2 \mathbb{1}_{]0,\theta]} (x)
$$

où $\theta > 0$ est un paramètre inconnu, et $C_\theta$ une constante ne dépendant que de $\theta$.

\begin{enumerate}
\item Justifier que $C_\theta = \frac{3}{\theta^3}$.
\item 
\begin{enumerate}
\item Calculer $\mathbb{E}_\theta [X_1]$
\item En déduire un estimateur $\hat{\theta}^{MM}_n$ par la méthode des moments.
\item Calculer le risque quadratique de $\hat{\theta}^{MM}_n$.
\item Donner la loi limite de $\hat{\theta}^{MM}_n - \theta$ après une renormalisation convenable.
\end{enumerate}
\item Calculer la fonction de répartition de $X_1$ notée $F_\theta$.
\item
\begin{enumerate}
\item Montrer que la médiane de $X_1$ est $x_{\frac{1}{2}} := F_\theta^{-1} (\frac{1}{2}) = 2^{-\frac{1}{3}} \theta$. En déduire une estimateur $\hat{\theta}^{m}_n$ consistant pour l'estimation de $\theta$.
\item Donner la loi limite de $\hat{\theta}^{m}_n - \theta$ après une renormalisation convenable.
\end{enumerate}
\item
\begin{enumerate}
\item Déterminer l'estimateur $\hat{\theta}^{MV}_n$ de $\theta$ par maximum de vraisemblance.
\item Déterminer la fonction de répartition de $v\hat{\theta}^{MV}_n$.
\item Montrer que $n(\theta-\hat{\theta}^{MV}_n)$ tend en loi vers une loi exponentielle dont on précisera le paramètre.
\end{enumerate}
\item Parmi les trois estimateurs, lequel faut-il privilégeir et porquoi?
\item Soit $\alpha \in ]0,1[$ Utiliser l'estimateur choisi pour construire un intervalle de confiance (asymptotique ou non) de niveau de condiance $1-\alpha$.
\end{enumerate}
\end{exercice}
\end{mdframed}

\end{document}